% Vorlage für eine Bachelorarbeit - 2012-2013 Timo Bingmann

% Dies ist nur eine Vorlage. Strikte Vorgaben wie die Bachelorarbeit auszusehen
% hat gibt es nicht. Darum können auch alle Teile angepasst werden.

\documentclass[12pt,a4paper,twoside]{scrartcl}

% Diese (und weitere) Eingabedateien sind in UTF-8
\usepackage[utf8]{inputenc}

% Verwende gute Type 1 Font: Latin Modern
\usepackage[T1]{fontenc}
\usepackage{lmodern}

% Sprache des Dokuments (für Silbentrennung und mehr)
\usepackage[german,english]{babel}

% Seitengröße - verwende fast die ganze A4 Seite
\usepackage[tmargin=22mm,bmargin=22mm,lmargin=20mm,rmargin=20mm]{geometry}

% Einrückung und Abstand zwischen Paragraphen
\setlength\parskip{\smallskipamount}
\setlength\parindent{0pt}

% Einige Standard-Mathematik Pakete
\usepackage{latexsym,amsmath,amssymb,mathtools,textcomp}

% Unterstützung für Sätze und Definitionen
\usepackage{amsthm}


\newtheorem{Satz}{Satz}[section]
\newtheorem{Definition}[Satz]{Definition}
\newtheorem{Lemma}[Satz]{Lemma}

\numberwithin{equation}{section}

% Deutsches Literaturverzeichnis
\usepackage{bibgerm}

% Unterstützung zum Einbinden von Graphiken
\usepackage{graphicx}

% Pakete die tabular und array verbessern
\usepackage{array,multirow}

% Kleiner enumerate und itemize Umgebungen
\usepackage{enumitem}

\setlist[enumerate]{topsep=0pt}
\setlist[itemize]{topsep=0pt}
\setlist[description]{font=\normalfont,topsep=0pt}

\setlist[enumerate,1]{label=(\roman*)}

% TikZ für Graphiken in LaTeX
\usepackage{tikz}
\usetikzlibrary{calc}

% Aktuelle Section und Untersection am Seitenkopf
\usepackage{fancyhdr}

\fancypagestyle{plain}{
  \fancyhead{}
  \fancyfoot{}
  \fancyfoot[LE,RO]{\normalsize\thepage}
  \renewcommand{\headrulewidth}{0pt}
  \renewcommand{\footrulewidth}{0pt}
}

\fancypagestyle{normal}{
  \setlength{\headheight}{20pt}
  \setlength\footskip{32pt}
  \fancyhead{}
  \fancyhead[LE]{\normalsize\textsc{\nouppercase{\leftmark}}}
  \fancyhead[RO]{\normalsize\textsc{\nouppercase{\rightmark}}}
  \fancyfoot{}
  \fancyfoot[LE,RO]{\normalsize\thepage}
  \renewcommand{\headrulewidth}{0.4pt}
  \renewcommand{\footrulewidth}{0pt}
}

% Hyperref für Hyperlink und Sprungtexte
\usepackage{xcolor,hyperref}

\hypersetup{
  pdftitle={High Quality Hypergraph Partitioning via Max-Flow-Min-Cut Computations},
  pdfauthor={Tobias Heuer},
  pdfsubject={Stichworte, weiteres Stichwort},
  colorlinks=true,
  pdfborder={0 0 0},
  bookmarksopen=true,
  bookmarksopenlevel=1,
  bookmarksnumbered=true,
  linkcolor=blue!60!black,
  %linkcolor=black,
  citecolor=blue!60!black,
  urlcolor=blue!60!black,
  filecolor=green!60!black,
  pdfpagemode=UseNone,
  unicode=true,
}

% Paket zum Setzen von Algorithmen in Pseudocode mit kleinen Stilanpassungen
\usepackage[ruled,vlined,linesnumbered,norelsize]{algorithm2e}
\DontPrintSemicolon
\def\NlSty#1{\textnormal{\fontsize{8}{10}\selectfont{}#1}}
\SetKwSty{texttt}
\SetCommentSty{emph}
\def\listalgorithmcfname{Algorithmenverzeichnis}
\def\algorithmautorefname{Algorithmus}
\let\chapter=\section % repariert ein Problem mit algorithm2e

\newtheorem{definition}{Definition}[section]
\newtheorem{sentence}{Sentence}[section]
\newtheorem{theorem}{Theorem}[section]
\newtheorem{lemma}{Lemma}[section]
\newtheorem{problem}{Problem}[section]

\newcommand{\BigO}[1]{$\mathcal O(#1)$}

\newcommand{\HypergraphDef}{$H = (V,E,c,\omega)$}
\newcommand{\FlowNetworkDef}{$G = (V',E',c')$}
\newcommand{\ShortT}[1]{{$T_{#1}(H)$}}
\newcommand{\ShortExT}[2]{{$T_{#1}(H,#2)$}}
\newcommand{\T}[1]{{$T_{#1}(H) = (V_{#1}, E_{#1},c_{#1})$}}
\newcommand{\ExtendedT}[2]{{$T_{#1}(H,#2) = (V_{#1}, E_{#1},c_{#1})$}}
\newcommand{\Hybrid}{{$T_{\text{Hybrid}}(H,V') = (V_{\text{Hybrid}}, E_{\text{Hybrid}},c_{\text{Hybrid}})$}}
\newcommand{\ShortHybrid}{{$T_{\text{Hybrid}}(H,V')$}}


%Residual Capacity
\newcommand{\R}[1]{{r_f({#1})}}
%Vertex Seperator
\newcommand{\V}[1]{V_S(#1)}

\newcommand{\DFS}{\emph{DFS}}
\newcommand{\BFS}{\emph{BFS}}

\newcommand{\GoldbergTarjan}{\textsc{GoldbergTarjan}}
\newcommand{\EdmondKarp}{\textsc{EdmondKarp}}

%Benchmark Instances
\newcommand{\ALL}{\sc All}
\newcommand{\Dual}{\sc Dual}
\newcommand{\Primal}{\sc Primal}
\newcommand{\ISPD}{\sc Ispd98}
\newcommand{\Literal}{\sc{Literal}}
\newcommand{\SPM}{\sc Spm}
\newcommand{\DAC}{\sc Dac}

%Experiment Flow Networks
\newcommand{\ExpLawler}{T_L}
\newcommand{\ExpNodeDegree}{T_H}
\newcommand{\ExpEdgeSize}{T_G}
\newcommand{\ExpHybrid}{T_{\text{Hybrid}}}

%Flow Variants
\newcommand{\FlowVariant}[4]{(#1\text{F},#2\text{C},#3\text{M},#4\text{FM})}

\newcounter{todocount}
\setcounter{todocount}{1}
\newcommand{\todo}[1]{{\color{red}\textbf{TODO \the\value{todocount}:}
        \emph{#1} \addtocounter{todocount}{1}}}

\begin{document}

%%%%%%%%%%%%%%%%%%%%%%%%%%%%%%%%%%%%%%%%%%%%%%%%%%%%%%%%%%%%%%%%%%%%%%

\pagestyle{empty} % keine Seitenzahlen

% Titelblatt der Arbeit
\begin{titlepage}

  \begin{center}\large
 
    \quad\includegraphics[height=17mm]{../img/logo/kit_logo_en.pdf} \hfill
    \includegraphics[height=20mm]{../img/logo/grouplogo-algo-blue.pdf}\quad\null

    \vfill

    Masterthesis
    \vspace*{2cm}

    {\bf\huge High Quality Hypergraph Partitioning \\ via Max-Flow-Min-Cut Computations \par}
    % Siehe auch oben die Felder pdftitle={}
    % mit \par am Ende stimmt der Zeilenabstand

    \vfill

    Tobias Heuer

    \vspace*{15mm}

    Abgabedatum: 19.10.2012

    \vspace*{45mm}

    \begin{tabular}{rl}
      Betreuer: & Prof. Dr. Peter Sanders \\
      & Sebastian Schlag\\
    \end{tabular}
    
    \vspace*{10mm}

    Institut für Theoretische Informatik, Algorithmik \\
    Fakultät für Informatik \\
    Karlsruher Institut für Technologie

    % English:
    % Institute of Theoretical Informatics, Algorithmics \\
    % Department of Informatics \\
    % Karlsruhe Institute of Technology

    \vspace*{12mm}
  \end{center}

\end{titlepage}

%%%%%%%%%%%%%%%%%%%%%%%%%%%%%%%%%%%%%%%%%%%%%%%%%%%%%%%%%%%%%%%%%%%%%%

\vspace*{0pt}\vfill

\hrule\medskip

Hiermit versichere ich, dass ich diese Arbeit selbständig verfasst und keine anderen, als die angegebenen Quellen und Hilfsmittel benutzt, die wörtlich oder inhaltlich übernommenen Stellen als solche kenntlich gemacht und die Satzung des Karlsruher Instituts für Technologie zur Sicherung guter wissenschaftlicher Praxis in der jeweils gültigen Fassung beachtet habe.

\bigskip

\noindent
Ort, den Datum

% Unterschrift (handgeschrieben)

\vspace*{5cm}

\clearpage

%%%%%%%%%%%%%%%%%%%%%%%%%%%%%%%%%%%%%%%%%%%%%%%%%%%%%%%%%%%%%%%%%%%%%%

\vspace*{0pt}\vfill

\selectlanguage{german}
\begin{abstract}
\centerline{\bf Zusammenfassung}

Hier die deutsche Zusammenfassung.

Ich bin Blindtext. Von Geburt an. Es hat lange gedauert, bis ich begriffen habe, was es bedeutet, ein blinder Text zu sein: Man macht keinen Sinn. Man wirkt hier und da aus dem Zusammenhang gerissen. Oft wird man gar nicht erst gelesen. Aber bin ich deshalb ein schlechter Text? Ich weiß, dass ich nie die Chance haben werde im Stern zu erscheinen. Aber bin ich darum weniger wichtig? Ich bin blind! Aber ich bin gerne Text. Und sollten Sie mich jetzt tatsächlich zu Ende lesen, dann habe ich etwas geschafft, was den meisten ,,normalen`` Texten nicht gelingt.

Ich bin Blindtext. Von Geburt an. Es hat lange gedauert, bis ich begriffen habe, was es bedeutet, ein blinder Text zu sein: Man macht keinen Sinn. Man wirkt hier und da aus dem Zusammenhang gerissen. Oft wird man gar nicht erst gelesen. Aber bin ich deshalb ein schlechter Text? Ich weiß, dass ich nie die Chance haben werde im Stern zu erscheinen. Aber bin ich darum weniger wichtig? Ich bin blind! Aber ich bin gerne Text.

\end{abstract}

\vfill

\selectlanguage{english}
\begin{abstract}
\centerline{\bf Abstract}

And here an English translation of the German abstract.

I'm blind text. From birth. It took a long time until I realized what it means to be random text: You make no sense. You stand here and there out of context. Frequently, they do not even read. But I have a bad copy? I know that I will never have the chance of appearing in the. But I'm any less important? I'm blind! But I like to text. And you should see me now actually over, then I have accomplished something that is not possible in most ``normal'' copies.

I'm blind text. From birth. It took a long time until I realized what it means to be random text: You make no sense. You stand here and there out of context. Frequently, they do not even read. But I have a bad copy? I know that I will never have the chance of appearing in the. But I'm any less important? I'm blind! But I like to text.

\end{abstract}

\vfill\vfill\vfill
\clearpage

%%%%%%%%%%%%%%%%%%%%%%%%%%%%%%%%%%%%%%%%%%%%%%%%%%%%%%%%%%%%%%%%%%%%%%

\vspace*{0pt}\vfill

\section*{Acknowledgements}

Lorem ipsum dolor sit amet, consectetuer adipiscing elit. Aenean commodo ligula eget dolor. Aenean massa. Cum sociis natoque penatibus et magnis dis parturient montes, nascetur ridiculus mus. Donec quam felis, ultricies nec, pellentesque eu, pretium quis, sem. Nulla consequat massa quis enim. Donec pede justo, fringilla vel, aliquet nec, vulputate eget, arcu. In enim justo, rhoncus ut, imperdiet a, venenatis vitae, justo.

\vfill\vfill\vfill
\clearpage

%%%%%%%%%%%%%%%%%%%%%%%%%%%%%%%%%%%%%%%%%%%%%%%%%%%%%%%%%%%%%%%%%%%%%%

\pagestyle{normal}
% markiere sections im Seitenkopf links und subsections rechts
\renewcommand\sectionmark[1]{\markboth{\thesection\quad\MakeUppercase{#1}}{\thesection\quad\MakeUppercase{#1}}}
\renewcommand\subsectionmark[1]{\markright{\thesubsection\quad\MakeUppercase{#1}}}

% Inhaltsverzeichnis
\tableofcontents

\clearpage

%%%%%%%%%%%%%%%%%%%%%%%%%%%%%%%%%%%%%%%%%%%%%%%%%%%%%%%%%%%%%%%%%%%%%%

\listoffigures
\listoftables
\listofalgorithms

\clearpage

%%%%%%%%%%%%%%%%%%%%%%%%%%%%%%%%%%%%%%%%%%%%%%%%%%%%%%%%%%%%%%%%%%%%%%

\section{Introduction}

\subsection{Motivation}

\subsection{Contributions}

\newpage

%%%%%%%%%%%%%%%%%%%%%%%%%%%%%%%%%%%%%%%%%%%%%%%%%%%%%%%%%%%%%%%%%%%%%%

\section{Preliminaries}

\subsection{Graphs}

\begin{definition}
A directed weighted graph $G = (V,E,c,\omega)$ is a set of nodes $V$ 
and a set of edges $E$ with a node weight function 
$c: V \rightarrow \mathbb{R}_{\ge 0}$ and an edge weight 
function $\omega: E \rightarrow \mathbb{R}_{\ge 0}$. A edge $e = (u,v)$ 
is a relation between two nodes $u,v \in V$.
\label{def:hypergraph}
\end{definition}

Two vertices $u$ and $v$ are \emph{adjacent}, if there exists an edge
$(u,v) \in E$. Two edges $e_1$ and $e_2$ are \emph{incident} to each other, if they
share a node. $I(v)$ denotes the set of all \emph{adjacent} nodes of 
$v$. The \emph{degree} of a node $v$ is $d(v) = |I(v)|$.

\begin{definition}
Given a directed graph $G = (V,E)$. A contraction of two nodes
$u$ and $v$ results in a new graph $G_{(u,v)} = (V\setminus\{v\},E')$, where
each edge of the form $(v,w)$ or $(w,v)$ in $E$ is replaced with an edge 
$(u,w)$ or $(w,u)$ in $E'$.
\label{def:contraction}
\end{definition}

A \emph{path} $P = (v_1,\ldots,v_k)$ is a sequence of nodes, where for
each $i \in [1,k-1]: (v_i,v_{i+1}) \in E$. A \emph{cycle} is a \emph{path}
$P = (v_1,\ldots,v_k)$ with $v_1 = v_k$. A \emph{strongly connected 
component} $C \subseteq V$ is a set of nodes where for each $u,v \in C$
exists a \emph{path} from $u$ to $v$. We can enumerate all \emph{strongly
connected components} (\emph{SCC}) in a directed graph $G$ with a linear time algorithm
proposed by Tarjan \cite{tarjan1972depth}. A directed graph $G$ without
any \emph{cycles} is called \emph{directed acyclic graph} (\emph{DAG}). On such
graphs we can define a \emph{topological order} $\gamma: V \rightarrow \mathbb{N}_+$ such
that for each $(u,v) \in E: \gamma(u) < \gamma(v)$. A \emph{topological order}
of a \emph{DAG} can be found in linear time with Kahn's algorithm \cite{kahn1962topological}.
We can transform a general directed graph $G$ into a \emph{DAG}, if we contract
each \emph{strongly connected component}. All concepts are illustrated in 
\autoref{img:scctoposort}.

\begin{figure}
\centering
\includegraphics[width=0.9\textwidth]{../img/preliminaries/scctoposort.eps}
\caption{Example of \emph{strongly connected components} of a directed graph and
         a \emph{topological order} on a \emph{directed acyclic graph}. Each 
         \emph{SCC} is marked in the same color.} 
\label{img:scctoposort}
\end{figure}

\begin{definition}
Let $G_{V'} = (V',E_{V'},c,\omega)$ be the subgraph of a graph $G$
induced by $V' \subseteq V$ with $E_{V'} = \{(u,v) \in E\ |\ u,v \in V'\}$.
\label{def:subgraph}
\end{definition}


\subsection{Flows and Applications}
\label{sec:applications}

Given a graph $G = (V,E,c)$ with capacity function $c: E \rightarrow \mathbb{R}_+$ and a source 
$s \in V$ and a sink $t \in V$. The maximum flow problem is about finding 
the maximum amount of flow from $s$ to $t$ in $G$. A flow is a function 
$f: E \rightarrow \mathbb{R}_+$, which have to satisfy the following constraints:
\begin{enumerate}
\item $\forall (u,v) \in E: f(u,v) \le c(u,v)$ (capacity constraint)
\item $\forall v \in V \setminus \{s,t\}: \sum_{(u,v) \in E} f(u,v) = \sum_{(v,u) \in E} f(v,u)$ (conservation of flow constraint)
\end{enumerate}
The capacity constraint restricts the flow on an edge $(u,v)$ by its capacity 
$c(u,v)$. Whereas the conservation of flow constraint ensures that the amount
of flow entering a node $v \in V \setminus \{s,t\}$ is the same as leaving a node.
The value of the flow is defined as $|f| = \sum_{(s,v) \in E} f(s,v) = \sum_{(v,t) \in E} f(v,t)$.
A flow $f$ is maximal, if there exists no other flow $f'$ with $|f'| > |f|$. \\
Another useful construct in connection with maximum flows, is the concept of the
\emph{residual graph} $G_f$ and the \emph{residual capacity} $r_f$ of a flow function $f$ on graph $G$.
The \emph{residual capacity} $r_f: V \times V \rightarrow \mathbb{R}_+$ is defined as follows:
\begin{enumerate}
\item $\forall (u,v) \in E: r_f(u,v) = c(u,v) - f(u,v)$
\item $\forall (u,v) \in E:$ If $f(u,v) > 0$ and $c(v,u) = 0$, then $r_f(v,u) = f(u,v)$
\end{enumerate}
For a edge $e = (u,v) \in E$ the residual capacity $r_f(u,v)$ is the remaining amount of 
flow which can be send over edge $e$. For each reverse edge $\overleftarrow{e} \notin E$ the
residual capacity $r_f(\overleftarrow{e})$ is the amount of flow which is send over $e$.
The \emph{residual graph} $G_f = (V,E_f,r_f)$ is the network containing all $(u,v) \in V \times V$
with $r_f(u,v) > 0$. More formal $E_f = \{(u,v)\ |\ r_f(u,v) > 0, (u,v) \in V \times V\}$.
\autoref{img:maximum_flow_example}~illustrates all presented concepts.

\begin{figure}
\centering
\includegraphics[width=0.9\textwidth]{../img/maximum_flow/maximum_flow_example.eps}
\caption{Figure illustrates concepts related to the maximum flow problem. A valid flow $f$ 
(red values) from $s$ to $t$ on a graph $G$ is shown on the left side. The corresponding
\emph{residual graph} $G_f$ with its \emph{residual capacities} (black values) 
is illustrated on the right side. The red highlighted path represents an \emph{augmenting path}
in $G$.}
\label{img:maximum_flow_example}
\end{figure}

The \emph{Max-Flow-Min-Cut}-Theorem is fundamental for many applications related to the maximum
flow problem \cite{ford1956maximal}.

\begin{theorem}
The value of a maximum $(s,t)$-flow obtainable in a graph $G$ is equal with the weight
of the minimum cutset in $G$ seperating $s$ and $t$.
\end{theorem}

Let $f$ be a maximum flow in a graph $G = (V,E,\omega)$ with $s \in V$ and $t \in V$. 
Further, let $A$ be the set containing all $v \in V$, which are \emph{reachable} from $s$
in $G_f$. A node $v$ is \emph{reachable} from a node $u$, if there exists a path from $u$
to $v$. Then the set of all cut edges between the bipartition $(A,V\setminus A)$ 
is a minimum-weight $(s,t)$-cutset \cite{ford2015flows}. $A$ can be calculated with a simple \BFS~in $G_f$ starting
from $s$. \\
From this analogy many solutions for related problems arose. Samples are listed below:
\begin{enumerate}
\item Maximum Bipartite-Matching
\item Minimum-Weight Vertex Seperator
\item Number of Edge-Disjoint Paths
\item Number of Vertex-Disjoint Paths
\end{enumerate}
Solutions for those problems sometimes involves a transformation $T$ of the graph $G$
into a flow network $T(G)$, such that the \emph{Max-Flow-Min-Cut}-Theorem is applicable. 
A problem important for this work is to find a minimum-weight $(s,t)$-vertex seperator
in a graph $G = (V,E,c)$ with $c: V \rightarrow \mathbb{R}_+$.

\begin{definition}
Let $G = (V,E,c)$ be a graph with $c: V \rightarrow \mathbb{R}_+$. $S \subseteq V$
is a vertex seperator for non-adjacent verticies $s \in V$ and $t \in V$ if the
removal of $S$ from graph $G$ seperates $s$ and $t$ ($s$ not \emph{reachable} from $t$).
A vertex seperator $S$ is a minimum-weight $(s,t)$-vertex seperator, if for all $S' \subseteq V$
$c(S) \le c(S')$.
\end{definition}

We can calculate a minimum-weight $(s,t)$-vertex seperator with a maximum flow
calculation in the following flow network (\todo{reference}):

\begin{definition}
\label{def:vertex_seperator_transformation}
Let $T_V$ be a transformation of a graph $G = (V,E,c)$ into 
a flow network $T_V(G) = (V_V, E_V, c_V)$ (with $c_V: E_V \rightarrow \mathbb{R}_+$). 
$T_V$ is defined as follows:
\begin{enumerate}
\item $V_V = \bigcup\limits_{v \in V}\ \{v', v''\}$
\item $\forall v \in V$ we add a directed edge $(v',v'')$
      with capacity $c_V(v',v'') = c(v)$
\item $\forall (u,v) \in E$ we add two directed edges $(u'', v')$ and 
      $(v'', u')$ with capacity $c_V(u'', v') = c_V(v'', u') = \infty$.
\end{enumerate} 
\end{definition}

The vertex seperator problem and transformation $T_V(G)$ are illustrated in \autoref{img:vertex_seperator_example}.
Obviously no edge between two adjacent nodes can be in a minimum-capacity $(s,t)$-cutset of $T_V(G)$,
because for all those edges the capacity is $\infty$. Therefore, the cutset must consist
of edges of the form $(v',v'')$. A minimum-weight $(s,t)$-vertex seperator can be calculated by
finding a maximum flow in $T_V(G)$, finding the minimum-capacity $(s,t)$-cutset with the procedure
described above and then map each cut edge $(v',v'')$ to their corresponding node $v$.\\
Given a set of sources $S$ and sinks $T$. The \emph{multi-source multi-sink} maximum flow problem is
about finding a maximum flow $f$ from all source nodes $s \in S$ to all sink nodes $t \in T$.
We can transform such a problem into a \emph{single-source single-sink} problem by adding
two additional nodes $s$ and $t$. We add a directed edge from $s$ to all source nodes $s' \in S$ 
and for all sink nodes $t' \in T$ a directed edge to $t$ with capacity 
$c(s,s') = c(t',t) = \infty$.

\begin{figure}
\centering
\includegraphics[width=0.9\textwidth]{../img/maximum_flow/vertex_seperator_example.eps}
\caption{ Illustration of the vertex seperator problem and the transformation $T_V(G)$ in which
          we can find a minimum vertex seperator with maximum flow computation. }
\label{img:vertex_seperator_example}
\end{figure}

\subsection{Hypergraphs}
\label{sec:hypergraph}

\begin{definition}
An undirected weighted hypergraph $H = (V,E,c,\omega)$ is a set of hypernodes $V$ 
and a set of hyperedges $E$ with a hypernode weight function 
$c: V \rightarrow \mathbb{R}_{\ge 0}$ and a hyperedge weight 
function $\omega: E \rightarrow \mathbb{R}_{\ge 0}$. A hyperedge $e$ 
is a subset of $V$ (formally: $\forall e \in E: e \subseteq V$).
\label{def:hypergraph}
\end{definition}

A hypergraph generalizes a graph by extending the definition of an edge, which 
can contain more than two nodes. Hyperedges are also called \emph{nets} and the hypernodes
of a net are called \emph{pins}. For a subset $V' \subseteq V$ and $E' \subseteq E$ we
define
\begin{align*}
c(V') = \sum_{v \in V'} c(v) \\
\omega(E') = \sum_{e \in E'} \omega(e)
\end{align*}

A vertex $v$ is \emph{incident} to a hyperedge $e$, if $v \in e$.
Two vertices $u$ and $v$ are \emph{adjacent}, if there exists an 
$e \in E$ such that $u,v \in e$. $I(v)$ denotes the set of all 
\emph{incident} nets of $v$. The \emph{degree} of a hypernode 
$v$ is $d(v) = |I(v)|$. The size of a net $e$ is the cardinality $|e|$.

\begin{definition}
Let $H_{V'} = (V',E_{V'},c,\omega)$ be the subhypergraph of a hypergraph $H$
induced by $V' \subseteq V$ with $E_{V'} = \{e \cap V'\ |\ e \in E: e 
\cap V' \neq \emptyset\}$.
\label{def:subhypergraph}
\end{definition}

A hypergraph $H = (V,E,c,\omega)$ can be represented as an undirected graph. 
There are two common transformations, called \emph{clique} and \emph{bipartite} 
representation \cite{HuMoerder85}. The \emph{clique} graph $G_x(H) = (V,E_x)$ models
each net $e$ as a clique between its pins. The \emph{bipartite} graph $G_*(H) = 
(V \cup E, E_*)$ contains all hypernodes and hyperedges as nodes and connects each
net $e$ with an undirected edge $\{e,v\}$ to all its pins $v \in e$. The two transformations
are illustrated in \autoref{img:hypergraph_transformation}.

\begin{figure}
\centering
\includegraphics[width=1.0\textwidth]{../img/preliminaries/hypergraph_transformation.eps}
\caption{Example of a hypergraph $H$ and its two corresponding graph representations.} 
\label{img:hypergraph_transformation}
\end{figure}

\subsection{Hypergraph Partitioning}

\begin{definition}
A $k$-way partition of a hypergraph $H$ is a partition of its hypernodes into
$k$ disjoint blocks $\Pi = \{V_1,\ldots,V_k\}$ such that $\bigcup_{i=1}^{k} V_i = V$
and $V_i \neq \emptyset$.
\label{def:kway_partition}
\end{definition}

For a $k$-way partition $\Pi = \{V_1,\ldots,V_k\}$, we define the \emph{connectivity set} of a
hyperedge $e$ with $\Lambda(e,\Pi) = \{V_i \in \Pi\ |\ V_i \cap e \neq \emptyset\}$. The \emph{connectivity}
of a net $e$ is $\lambda(e,\Pi) = |\Lambda(e,\Pi)|$. A hyperedge $e$ is \emph{cut}, if
$\lambda(e,\Pi) > 1$. $E(\Pi) = \{e\ |\ \lambda(e,\Pi) > 1\}$ is the set of all \emph{cut} 
nets. We say two blocks $V_i$ and $V_j$ are adjacent, if there exists a hyperedge
$e$ with $V_i,V_j \in \Lambda(e,\Pi)$.

\begin{definition}
For a $k$-way partition $\Pi = \{V_1,\ldots,V_k\}$ of a hypergraph $H$ 
the quotient graph $Q = (\Pi,E')$ is a undirected graph containing an 
edge between each pair of adjacent blocks of $\Pi$.
More formal, $E' = \{(V_i,V_j)\ |\ \exists e \in E: V_i,V_j \in \Lambda(e,\Pi)\}$
\label{def:quotient_graph}
\end{definition}

We say a $k$-way partition is $\epsilon$-balanced, if each block 
$V_i \in \Pi$ satisfies the \emph{balance constraint} 
$c(V_i) \le (1+\epsilon)\lceil\frac{c(V)}{k}\rceil$.

\begin{definition}
The $k$-way hypergraph partitioning problem is to find an $\epsilon$-balanced $k$-way
partition $\Pi$ of a hypergraph $H$ such that a certain objective function is minimized.
\label{def:kway_partitioning_problem}
\end{definition}

There exists several objective function in the hypergraph partitioning context,
which should either be minimized or maximized. The most popular objective function 
is the $\text{cut}$ metric (especially for \emph{graph partitioning}), which is defined as
\[\omega_H(\Pi) = \sum_{e \in E(\Pi)} \omega(e)\]
The goal is to minimize the sum of all \emph{cut} hyperedges. Another important metric
for this work is the $(\lambda - 1)$-metric or \emph{connectivity} metric, 
which is defined as
\[(\lambda - 1)_H(\Pi) = \sum_{e \in E} (\lambda(e) - 1)\omega(e)\]
The idea behind this function is to minimize the \emph{connectivity} of all hyperedges.

\newpage

%%%%%%%%%%%%%%%%%%%%%%%%%%%%%%%%%%%%%%%%%%%%%%%%%%%%%%%%%%%%%%%%%%%%%%

\section{Related Work}

%%%%%%%%%%%%%%%%%%%%%%%%%%%%%%%%%%%%%%%%%%%%%%%%%%%%%%%%%%%%%%%%%%%%%%

\subsection{Maximum Flow Algorithms}
\label{sec:max_flow}

In Section \ref{sec:applications} we introduce the concept of flows in a network. We
will now present two algorithms to solve the maximum flow problem.

\subsubsection{Augmenting-Path Algorithms}
\label{sec:aug_path}

An \emph{augmenting path} $P = \{v_1,\ldots,v_k\}$ is a path in $G_f$ with $v_1 = s$ and 
$v_k = t$ \cite{edmonds1972theoretical}. \autoref{img:edmond_karp_example}~illustrates such a path.
Since all $(v_i,v_{i+1}) \in G_f$ it follows that $r_f(v_i,v_{i+1}) > 0$. 
Therefore, we can increase the flow on all edges $(v_i,v_{i+1})$ by 
$\Delta f = \min_{i \in [1,\ldots,k-1]} r_f(v_i,v_{i+1})$. It can be shown that $f$ is not a
maximum flow if an \emph{augmenting path} exists in $G_f$ \cite{edmonds1972theoretical}. \\
One way to calculate a maximum flow $f$ is to find \emph{augmenting paths} in $G_f$ as
long as there exists one. The algorithm was established by Ford and Fulkerson \cite{ford1956maximal} and
consists of two phases. First, we search for an \emph{augmenting path} $P = \{v_1,\ldots,v_k\}$
from $s$ to $t$, e.g., with a simple \DFS. Afterwards, we increase the flow on each
edge $(v_i,v_{i+1})$ by $\Delta f$ and decrease the flow on each reverse edge $(v_{i+1},v_i)$
by $\Delta f$. If the capacities are integral, the algorithm always terminates. Since we can find an \emph{augmenting
path} in $G_f$ with a simple \DFS~in \BigO{|V|+|E|} and increase the
flow on every path by at least one, the running time of the algorithm can be bounded by \BigO{|E||f_{max}|}.
We can construct instances, where the running time is \BigO{|E||f_{max}|} or even the maximum flow $|f_{max}|$ 
is exponential in the problem size \cite{edmonds1972theoretical}. \\
Edmond and Karp \cite{edmonds1972theoretical} improved Ford \& Fulkerson's algorithm by increasing the flow
along an \emph{augmenting path} of minimal length. The shortest path from $s$ to $t$ in a 
graph with unit lengths can be found by a simple \BFS~calculation. It can be shown, that the
total number of \emph{augmentations} is \BigO{|V||E|}. The running time of Edmond \& Karp's
maximum flow algorithm is \BigO{|V||E|^2}. A sample execution of the algorithm
is presented in \autoref{img:edmond_karp_example}.

\begin{figure}
\centering
\includegraphics[width=0.9\textwidth]{../img/maximum_flow/edmond_karp_example.eps}
\caption{Execution of Edmond \& Karps maximum flow algorithm \cite{edmonds1972theoretical}.
         The network $G$ with its capacities $c$ (black values) and flow $f$ (red values) is illustrated
         on the left side. The residual graph $G_f$ with its \emph{residual capacities} $r_f$ (black values)
         is presented on the right side. In each step the current \emph{augmenting path} in $G_f$ is highlighted
         by a red path. }
\label{img:edmond_karp_example}
\end{figure}

\subsubsection{Push-Relabel Algorithm}
\label{sec:push_relabel}

Goldberg and Tarjan \cite{goldberg1988new} implemented a maximum flow algorithm
not based on finding an \emph{augmenting path} in the \emph{residual graph}. The idea is
to maintain a \emph{preflow} during the execution of the algorithm which satisfies the capacity 
constraints, but only a weakened form of the conservation of flow constraint:
\[\forall v \in V \setminus \{s,t\}: \sum_{u \in V} f(v,u) \le \sum_{u \in V} f(u,v)\]
The algorithm maintains a \emph{distance labeling} $d: V \rightarrow \mathbb{N}$ and an 
\emph{excess function} $e_f: V \rightarrow \mathbb{N}$. The \emph{distance labeling} satisfies
the following conditions: $d(s) = |V|$, $d(t) = 0$ and for each $(u,v) \in E_f$, $d(u) \le d(v) + 1$. We say an
residual edge $(u,v)$ is \emph{admissible} if $d(u) = d(v) + 1$. A node $v$ is \emph{active}
if $v \notin \{s,t\}$ and $e_f(v) > 0$.\\
Initially, all \emph{labels} and \emph{excess} values are set to zero except source node $s$ will be set to $d(s) = 1$
and $e_f(s) = \infty$. For each \emph{active} node $u$ the algorithm performs two update
operations, called \emph{push} and \emph{relabel}. The first operation pushes flow
over each \emph{admissible} edge $(u,v)$. After a \emph{push} $e_f(u) = e_f(u) - 
\min{(e_f(u),r_f(u,v))}$ and $e_f(v) = e_f(v) + \min{(e_f(u),r_f(u,v))}$. If there is no
\emph{admissible} edge, a \emph{relabel} operation is performed, which replaces $d(u)$ by
$\min_{(u,v) \in E_f} d(v) + 1$. The algorithm terminates, if none of the nodes is \emph{active}.
The worst case complexity of the algorithm is \BigO{n^3}. The running time can be reduced
to \BigO{n^2\log{n}} with \emph{Dynamic Trees} \cite{goldberg1988new, sleator1981data}, but this
implementation is not practical due to a large hidden constant factor.\\
The \emph{push-relabel} algorithm is one of the fastest maximum flow algorithms in practice
because there exist several speed-up techniques. The first one is
the \emph{global relabeling} heuristic which frequently updates the \emph{distance labels} by computing
the shortest path in the residual graph from all nodes to the sink \cite{cherkassky1997implementing}.
This can be done with a backward \emph{BFS} in linear time. This heuristic is performed periodically,
e.g., after every $n$ relabeling. \\
The second heuristic is the \emph{gap heuristic} \cite{cherkassky1994fast,derigs1989implementing}.
If at a particular stage of the algorithm there is no node $u$ with $d(u) = g < n$, then for each node
$v$ with $g < d(v) < n$ the sink is not reachable anymore. Therefore, we can increase the \emph{distance
label} of all those nodes to $n$. To implement this heuristic, we maintain a linked list of
nodes with distance label $i$.

%%%%%%%%%%%%%%%%%%%%%%%%%%%%%%%%%%%%%%%%%%%%%%%%%%%%%%%%%%%%%%%%%%%%%%

\subsection{Modeling Flows on Hypergraphs}
\label{sec:related_lawler}

Consider the \emph{bipartite graph} representation $G_*(H)$ of a hypergraph $H$ (see Section \ref{sec:hypergraph}).
Hu and Moerder \cite{HuMoerder85} introduce node capacities in $G_*(H)$. Each hyperedge
$e$ has a capacity equal to $\omega(e)$ and each hypernode has infinite capacity. 
Further, they show that a minimum-weight $(s,t)$-vertex separator in $G_*(H)$
is equal with a minimum-weight $(s,t)$-cutset of a hypergraph $H$. 
Finding such a separator is a flow problem and can be calculated with the flow network \ShortT{L} 
presented by Lawler \cite{lawler1973}:

\begin{definition}
Let $T_L$ be the transformation of a hypergraph \HypergraphDef~into 
a flow network \T{L} proposed by Lawler \cite{lawler1973}. \ShortT{L} is defined as follows:
\begin{enumerate}
\item $V_L = V \cup \bigcup\limits_{e \in E}\ \{\incoming{e}, \outgoing{e}\}$
\item $\forall e \in E$ we add a directed edge $(\incoming{e},\outgoing{e})$ 
      with capacity $\capa_H(\incoming{e},\outgoing{e}) = \omega(e)$
\item $\forall v \in V$ and $\forall e \in I(v)$ we add two directed edges $(v, \incoming{e})$ and 
      $(\outgoing{e}, v)$ with capacity $\capa_L(v, \incoming{e}) = \capa_L(\outgoing{e},v) = \infty$.
\end{enumerate} 
\end{definition}
\begin{figure}
\centering
\includegraphics[width=0.8\textwidth]{../img/network_transformation/lawler_transformation.eps}
\caption{Transformation of a hypergraph into a equivalent flow network by Lawler \cite{lawler1973}. Note,
capacity of the black edges in the flow network is $\infty$.}
\label{img:lawler_transformation}
\end{figure}
An example of this transformation is shown in \autoref{img:lawler_transformation}.
\ShortT{L} is nearly equivalent to the transformation $T_V(G)$ described in Defintion \autoref{def:vertex_seperator_transformation}
except that we do not have to split the hypernodes $v \in V$.
A hypernode cannot be in a minimum-capacity $(s,t)$-vertex separator, because each $v \in V$ has
infinity capacity \cite{HuMoerder85}. Therefore, a minimum-capacity $(s,t)$-cutset 
of \ShortT{L} is equal to a minimum $(s,t)$-vertex separator of $G_*(H)$.
The resulting graph \ShortT{L} has $|V_L| = 2|V| + |E|$ nodes and $|E_L| = 2(\bar{e}+1)|E|$ edges, where
$\bar{e}$ is the average size of a hyperedge \cite{pistorius2003}. Using \emph{Edmond-Karps}
maximum flow algorithm (see Section \ref{sec:aug_path}) on flow network \ShortT{L} 
takes time \BigO{|V|^2|E|^2} \cite{lawler1973}. \\
A minimum-weight $(s,t)$-cutset of $H$ can be found by simply mapping the minimum-capacity
$(s,t)$-cutset to their corresponding hyperedges in $H$ (see Section \ref{sec:applications}). 
The minimum-weight $(s,t)$-bipartition are all $v \in V$ \emph{reachable} from $s$ in the 
\emph{residual graph} of \ShortT{L} and the counterpart are all hypernodes not \emph{reachable}
from $s$. \\ 
In \autoref{img:lawler_augmenting_example} we illustrate the structure of \ShortT{L} and demonstrate
what happens after we \emph{augment} along a path in the Lawler-Network. This figure can be used as
a reference if you need an illustration of techniques used in the proofs of Section \ref{sec:opt_flow_network}.\\
In this thesis, we often have to mix up nodes and edges of $H$ and \ShortT{L}. If we use
$v \in V_L$, there also exists a corresponding $v \in V$. $v$ can be used in both contexts.
For all $e \in E$ there exists two corresponding nodes $\incoming{e},\outgoing{e} \in V_L$. $\incoming{e}$ 
is called \emph{incoming hyperedge node} and $\outgoing{e}$ is called \emph{outgoing hyperedge node}. 
In some cases we need to treat $\incoming{e},\outgoing{e} \in V_L$ the same way as their corresponding 
hyperedge $e \in E \Rightarrow \incoming{e_1} \cap \outgoing{e_2}$ or $\outgoing{e_1} \cap \incoming{e_2}$ should be the same as $e_1 \cap e_2$.
\begin{figure}[t!]
\centering
\includegraphics[width=0.75\textwidth]{../img/network_transformation/lawler_augmenting_example.eps}
\caption{Illustration of the structure of a hyperedge $e$ in \ShortT{L} and the effect of augmenting
         along a path in this network. The labeling on an edge represents the flow $f(x,y)$ and
         the capacity $\capacity{x,y}$ denoted with $f(x,y)/\capacity{x,y}$. The labeling of the
         edges in the residual graph denotes the residual capacity $r_f(x,y)$. The red highlighted
         path represents an augmenting path.}
\label{img:lawler_augmenting_example}
\end{figure}

%%%%%%%%%%%%%%%%%%%%%%%%%%%%%%%%%%%%%%%%%%%%%%%%%%%%%%%%%%%%%%%%%%%%%%

\subsection{Flow-based Local Search on Graphs}
\label{sec:flow_local_search_graph}

It seems naturals to utilize maximum flow computations to improve the cut metric of a 
given partition of a graph. Lang and Rao \cite{lang2004flow} use an approach,
called \emph{Max-Flow Quotient-cut Improvement} (MQI), to improve the quality
of a graph when metrics such as \emph{expansion} or \emph{conductance}
are used. For a given bipartition $(S,\bar{S})$, they find the best 
improvement among all bipartitions $(S',\bar{S'})$ such that $S' \subset S$
by solving a flow problem. Andersen and Lang \cite{andersen2008algorithm}
proposed a flow-based improvement algorithm, called \emph{Improve},
which works similar as MQI, but did not restrict the output of the 
partition to $S' \subset S$. However, both techniques can not guarantee 
that the resulting bipartition is balanced and only are applicable for $k=2$. \\
Schulz and Sanders \cite{sanders2011engineering} integrate flow-based refinement algorithm 
in their \emph{multilevel graph partitioner} \emph{KaFFPa}. In general,
they build a flow problem around a region $B$ of the cut and connect the \emph{border} 
of $B$ with the source resp. sink. $B$ is defined in such a way that the flow computation
yields a feasible cut in the original graph. Many ideas of this work are used in this
thesis and adapted to hypergraphs. Therefore, we will give a detailed description
of the concepts and advanced techniques to improve graph partitions.

\subsubsection{Balanced Bipartitioning}
\label{sec:balanced_bipartitioning}
\begin{figure}
\centering
\includegraphics[width=0.6\textwidth]{../img/flow_local_search/balanced_bipartitioning.eps}
\caption{Configuration of a flow problem around the cut of graph $G$ \cite{andersen2008algorithm}.}
\label{img:balanced_bipartition}
\end{figure}
Let $(V_1,V_2)$ be a balanced bipartition of a graph $G = (V,E,c,\omega)$. 
Further, $P(v) = 1$, if $v \in V_1$ and
$P(v) = 2$, otherwise. We will now explain how a given bipartition 
can be improved with flow computations. This technique can also be applied on a $k$-way 
partition by applying the approach on two adjacent blocks \cite{sanders2011engineering}. \\
Let $\delta := \{ u\ |\ \exists (u,v) \in E: P(u) \neq P(v) \}$ be the set of nodes
around the cut of $G$. For a set $B \subseteq V$ we define its border 
$\delta B := \{u \in B\ |\ \exists (u,v) \in E: v \notin B\}$.
The basic idea is to build a region $B$ around all cut nodes $\delta$ of 
$G$ and connect all nodes in $\delta B \cap V_1$ to the source node $s$ and all nodes in 
$\delta B \cap V_2$ to the sink node $t$. \\
We can construct $B := B_1 \cup B_2$ with two \emph{Breadth First Searches} (\BFS). 
One is initialized with all nodes $\delta \cap V_1$ and stops if $c(B_1)$ would 
exceed $(1+\epsilon)\frac{c(V)}{2} - c(V_2)$. The second is initialized with 
all nodes $\delta \cap V_2$ and stops if $c(B_2)$ would exceed 
$(1+\epsilon)\frac{c(V)}{2} - c(V_1)$. The two \BFS s only touch nodes of $V_1$ resp. $V_2 \Rightarrow$
$B_1 \subseteq V_1$ and $B_2 \subseteq V_2$. The constraints for the weights of $B_1$
and $B_2$ guarantes that the bipartition is still balanced after a \emph{Max-Flow-Min-Cut}
computation. Connecting $s$ resp. $t$ to all border nodes $\delta B \cap V_1$ resp.
$\delta B \cap V_2$ ensures that a non-cut edge not contained in $G_B$ is not a cut edge after
assigning the minimum $(s,t)$-bipartition of subgraph $G_B$ to $G$. This also yields 
the conclusion that each minimum $(s,t)$-cutset in $G_B$ leads to a cut smaller or 
equal to the old cut of $G$. All concepts are illustrated in \autoref{img:balanced_bipartition}.


\subsubsection{Adaptive Flow Iterations}
\label{sec:adaptive_flow_iterations}

Sanders and Schulz \cite{sanders2011engineering} introduce several techniques to improve
their basic approach. If the \emph{Max-Flow-Min-Cut} computation on $G_B$ leads to an
improved cut, we can apply the method described in Section 
\ref{sec:balanced_bipartitioning} again. An extension of this approach is to iteratively adapt the
size of the flow problem based on the result of the maximum flow computation.
We define $\epsilon' := \alpha\epsilon$ for a $\alpha \ge 1$ and let the size of $B$ depend
on $\epsilon'$ rather than on $\epsilon$. If we find an improvement on $G$, we
increase $\alpha$ to $\min\{2\alpha, \alpha'\}$ where $\alpha'$ is a predefined upper bound
for $\alpha$. If not, we decrease the size of $\alpha$ to 
$\max\{\frac{\alpha}{2},1\}$. This approach is called
\emph{adaptive flow iterations} \cite{sanders2011engineering}.

\subsubsection{Most Balanced Minimum Cut}
\label{sec:related_mbmc}

Picard and Queyranne \cite{picard1980structure} show that all minimum $(s,t)$-cutsets 
are computable with one maximum $(s,t)$-flow computation.
To understand the main theorem and the algorithm to compute all minimum $(s,t)$-cutsets we
need the definition of a \emph{closed node set} $C \subseteq V$ of a graph $G$.

\begin{definition}
Let $G = (V,E)$ be a graph and $C \subseteq V$. $C$ is called a closed node set iff the 
condition $u \in C$ implies that for all edges $(u,v) \in E$ also $v \in C$.
\end{definition}

A \emph{closed node set} is illustrated in \autoref{img:closed_node_set}. A simple observation
is that all nodes on a cycle have to be in the same \emph{closed node set} per definition. Therefore
we can contract all \emph{Strongly Connected Components} (SCC) of $G$ with a linear time algorithm
proposed by Tarjan \cite{tarjan1972depth} and sweep in reverse
topological order over the contracted graph to enumerate all \emph{closed node sets}. Note, if
we contract all SCCs of $G$ the resulting graph is a \emph{Directed Acyclic Graph} (DAC). Therefore, 
a topological order exists.\\
With the Theorem of Picard and Queyranne \cite{picard1980structure} we can enumerate
all minimum $(s,t)$-cuts of $G$ with one maximum flow computation.

\begin{theorem}
\label{theorem:mbmc}
There is a $1$-$1$ correspondence between the minimum $(s,t)$-cuts of a graph and the closed node
sets containing $s$ in the residual graph of a maximum $(s,t)$-flow.
\end{theorem}

All \emph{closed node sets} in the residual graph of $G$ induce a minimum $(s,t)$-cutset on $G$.
They can be calculated with the algorithm described above having the residual graph of
$G$ as input. The running time of the algorithm is \BigO{|V| + |E|}. \\
A common problem of the \emph{adaptive flow iteration} approach (see Section \ref{sec:adaptive_flow_iterations}) is
that using a large $\alpha$ often leads to cuts in $G$ that violate the balanced constraint. 
We can enumerate all minimum $(s,t)$-cutsets with one maximum 
flow computation and therefore have a higher probability to find
a feasible partition after a \emph{Max-Flow-Min-Cut} computation. We refer to this method as
\emph{Most Balanced Minimum Cut}.
\begin{figure}
\centering
\includegraphics[width=0.9\textwidth]{../img/flow_local_search/closed_node_set.eps}
\caption{$C = \{s,a,b,c\}$ is a \emph{closed node set} of graph $G$ (left side).
         After contracting all \emph{Strongly Connected Components}, we can enumerate all
         \emph{closed node sets} of $G$ by sweeping in reverse topological order over the
         contracted graph (right side).}
\label{img:closed_node_set}
\end{figure}


\subsubsection{Active Block Scheduling}
\label{sec:abs}
\emph{Active Block Scheduling} is a \emph{quotient graph style refinement} technique for
$k$-way partitions \cite{holtgrewe2010engineering,sanders2011engineering}. 
The algorithm is organized in rounds and executes a two-way 
local improvement algorithm on each adjacent pair of 
blocks in the \emph{quotient graph} where at least one of both is \emph{active}. 
Initially all blocks are \emph{active}. A block becomes \emph{inactive}
if its none of its nodes move in a round. The algorithm
terminates, if all blocks are \emph{inactive}. \\
Fiduccia and Mattheyses \cite{fiduccia1988linear} introduce a linear time
two-way local search heuristic, called \emph{FM} heuristic, 
which is fundamental for many graph partitioning algorithms.
They define the gain $g(v)$ of a node $v \in V$ as the reduction of the cut metric when
moving $v$ from its current block to an other block. By maintaining the gains of the
nodes in a special data structure, called \emph{bucket queue}, they can find a maximum
gain node in constant time. After moving a maximum gain node, they are also able to update the
data structure in time equal to the number of adjacent nodes.\\
The local improvement algorithm (for \emph{Active Block Scheduling}) can either 
be an \emph{FM} local search or a flow-based approach or even a combination of 
both as proposed by Sanders and Schulz \cite{sanders2011engineering}.

%%%%%%%%%%%%%%%%%%%%%%%%%%%%%%%%%%%%%%%%%%%%%%%%%%%%%%%%%%%%%%%%%%%%%%

\subsection{Hypergraph Partitioning}

In this Section, we review how most hypergraph partitioners solve the \emph{hypergraph
partitioning problem}. The most successful
approach is the \emph{multilevel paradigm} \cite{alpert1995recent,bader2013graph,
papa2007hypergraph} which we describe in Section \ref{sec:multilevel_paradigm}.
The results of this thesis is integrated into $n$-level hypergraph partitioner \emph{KaHyPar}. Therefore, we
give a brief overview of implementation details of this framework (see Section \ref{sec:kahypar}).

\subsubsection{Multilevel Paradigm}
\label{sec:multilevel_paradigm}

The \emph{multilevel paradigm} is a three phase algorithm to solve the \emph{hypergraph 
partitioning problem} (see \autoref{img:multilevel_paradigm}). In the first stage, called
\emph{coarsening phase}, vertex matchings or clusterings are calculated to be contracted. This process is
repeated until a predefined number of hypernodes remains. The sequence of successively
smaller hypergraphs is called \emph{levels}. If the hypergraph $H$ is small enough, we can use
expensive algorithms to \emph{initially partition} $H$ into $k$ blocks. Afterwards, we
\emph{uncontract} each \emph{level} in reverse order of \emph{contraction} by projecting
the partition to the next \emph{level}. After \emph{uncontraction} a \emph{refinement} heuristic can be
used to improve the quality of the current partition according
to an objective function. The most commonly used \emph{refinement} algorithm is the \emph{FM}
algorithm \cite{fiduccia1988linear}.

\begin{figure}
\centering
\includegraphics[width=1.0\textwidth]{../img/hypergraph_partitioning/multilevel.eps}
\caption{Multilevel Hypergraph Partitioning}
\label{img:multilevel_paradigm}
\end{figure}


\subsubsection{$n$-Level Hypergraph Partitioning}
\label{sec:kahypar}

\emph{KaHyPar} is a multilevel hypergraph partitioner in its most extreme version, which
removes only a single vertex in one \emph{level} of the hierarchy. It seems to be the method
of choice for optimizing cut- and the $(\lambda - 1)$-metric unless speed is more important than
quality \cite{heuer2017improving}. The framework provides a \emph{direct $k$-way} \cite{akhremtsev2017engineering} 
and a \emph{recursive bisection} mode, which recursively calculates bipartitions 
(with \emph{multilevel paradigm}) until the hypergraph is divided into $k$ blocks 
\cite{schlag2016k}. \emph{KaHyPar} consists of four phases: \emph{Preprocessing}
and the three stages of the \emph{multilevel paradigm}. \\
In the \emph{preprocessing} step community structures of the hypergraph are detected. The
hypergraph is transformed into a bipartite graph $G_*(H)$ (see Section \ref{sec:hypergraph}) and
a community detection algorithm is executed which optimizes \emph{modularity} \cite{fortunato2010community,
heuer2017improving}. During the \emph{coarsening phase} contractions are restricted to vertices within the same 
community. The contraction partners are chosen according to the \emph{heavy-edge} rating function
$r(u,v) := \sum_{e \in I(u) \cap I(v)} \frac{\omega(e)}{|e|-1}$ \cite{karypis1999multilevel}. The
function prefers vertices which share a large number of heavy nets with small size. The contraction
algorithm works in passes. At the beginning of each pass a random permutation of the vertices 
is created and for each vertex $u$, the contraction partner $v$ is determined according
to the \emph{heavy-edge} rating function \cite{schlag2016k}. A pass ends if each vertex is either
considered as representative or contraction partner. The passes are repeated until only
$t = 160k$ hypernodes remains.
The \emph{initial partitioning} uses the \emph{recursive bisection} approach to calculate
a $k$-way partition in combination with a portfolio of initial partitioning techniques 
\cite{heuer2015engineering}. In the \emph{refinement phase}, a localized \emph{FM} search is started \cite{fiduccia1988linear},
initialized with the current uncontracted vertices. The \emph{local search} maintains $k$ \emph{priority
queues} (PQ) for each block $V_i$ \cite{akhremtsev2017engineering}. A hypernode $v$ contained in 
the $i$-th PQ with gain $g$ means that moving vertex $v$ to block $V_i$ has gain $g$. After a move,
the gains of all adjacent hypernodes are updated with a \emph{delta-gain} update strategy \cite{papa2007hypergraph}.
The recalculation of all gain values at the beginning of a \emph{FM} pass is one of the main bottlenecks
of the algorithm \cite{papa2007hypergraph}. Therefore, Schlag et al. \cite{akhremtsev2017engineering,
schlag2016k} introduce a \emph{gain cache}, which prevents
expensive recalculations of the corresponding gain function. The \emph{gain cache} is maintained
with \emph{delta-gain} updates in the same way as the \emph{PQs}. Further, the \emph{local search}
is stopped, when an improvement during an \emph{FM} pass becomes unlikely. This model is
called \emph{adaptive stopping rule} \cite{akhremtsev2017engineering}. Sanders and
Osipov \cite{osipov2010n} shows that it is unlikely that \emph{local search} gives an improvement if
$p > \frac{\sigma^2}{4\mu^2}$, where $p$ is the number of moves in the current \emph{FM} pass,
$\mu$ is the average gain, and $\sigma^2$ the corresponding variance.

\newpage

%%%%%%%%%%%%%%%%%%%%%%%%%%%%%%%%%%%%%%%%%%%%%%%%%%%%%%%%%%%%%%%%%%%%%%

\section{Optimized Approach on Modelling Flows in Hypergraphs}
\label{sec:opt_flow_network}

In Section \ref{sec:related_lawler} we have shown how a hypergraph $H$ can be transformed into 
a flow network \ShortT{L} such that every minimum-weight $(S,T)$-cutset in $H$ is a
minimum-capacity $(S,T)$-cutset in \ShortT{L} \cite{lawler1973}. However, the resulting flow
network has significantly more nodes and edges than the original hypergraph. The running time
of a maximum $(S,T)$-flow algorithm depends heavily on the problem size. 
Therefore, different modeling approaches, which reduce the number of nodes and edges,
can have a crucial impact on the running time of the flow algorithm. \\
We will present techniques to sparsify the flow network 
proposed by Lawler. First, we will show how any subset $V' \subseteq V$ of hypernodes could be removed 
from \ShortT{L} (see Section \ref{sec:heuer_network}). This approach minimizes
the number of nodes, but in some cases, the number of edges can be
significantly higher than in \ShortT{L}. The basic idea of this technique 
can still be applied to remove low degree hypernodes from the Lawler-Network \emph{without} 
increasing the number of edges (see Section \ref{sec:degree_network}). Additionally, we show
how every hyperedge $e$ of size $2$ can be removed by inserting an undirected flow edge between
the corresponding nodes $v_1,v_2 \in e$  (see Section \ref{sec:edge_size_network}). 
Finally, we combine the two suggested approaches into a Hybrid-Network 
(see Section \ref{sec:hybrid_network}).



\subsection{Removing Hypernodes via Clique-Expansion}
\label{sec:heuer_network}

In this Section, we show how all hypernodes of \ShortT{L} can be removed such that a maximum 
$(S,T)$-flow on the new network induce a minimum-weight $(S,T)$-cutset on $H$. If a hypernode $v \in V$
occurs in an augmenting path $P$ the previous node in the path must be a hyperedge, either
$\incoming{e}$ or $\outgoing{e}$. Further, for all $e \in I(v)$ the capacity $\capa_L(v,\incoming{e})$ is $\infty$. Therefore, 
if we push flow over a hypernode $v$, coming from a hyperedge, we can redirect
the flow to any hyperedge node $\incoming{e} \in I(v)$ during the whole maximum flow calculation, because 
$\capa_L(v,\incoming{e}) = \infty$. The following lemma is central to our first sparsifying 
technique and is illusrated in \autoref{img:clique_expansion}. Given a graph $G = (V,E)$ we define the two sets
$in(u) := \{v\ |\ (v,u) \in E\}$ and $out(u) := \{v\ |\ (u,v) \in E\}$ with $u \in V$. 

\begin{lemma}
\label{lemma:node_removal}
Let $G = (V,E,\capa)$ be a flow network and $u \in V$ a node where
$\forall v \in in(u): \capacity{v,u} = \infty$ and $\forall w \in out(u): \capacity{u,w} = \infty$.
Further, let $G(u) = (V\setminus\{u\}, E_u, \capa_u)$ be the flow network obtained by removing
$u$ and inserting a shortcut edge between each $v \in in(u)$ and $w \in out(u)$ with $\capa_u(v,w) = \infty$.
Let $f$ be a maximum $(S,T)$-flow of G and $f'$ a maximum $(S,T)$-flow of
$G(u)$ if $|f| < \infty$ and $u \notin S\cup T$, then $|f| = |f'|$.
\end{lemma}

\begin{proof}
Let $f$ be a maximum $(S,T)$-flow of $G$. We define a maximum $(S,T)$-flow $f'$ 
of $G(u)$ as follows:
\begin{equation}
\label{eq:eq1}
f'(v,w) =  
  \begin{cases}
      \frac{f(v,u)f(u,w)}{\sum_{w \in out(u)} f(u,w)}, & \text{if } v \in in(u), w \in out(u) \\
      f(v,w), & \text{otherwise}
   \end{cases} 
\end{equation}
$f'$ is choosen in such a way that for all $v \in in(u): \sum_{w \in out(u)} f'(v,w) = f(v,u)$
and for all $w \in out(u): \sum_{v \in in(u)} f'(v,w) = f(u,w)$. Therefore, $f'$
satisfies the flow conservation constraint and is an valid flow function. Since
$u \notin S\cup T$, it follows that $|f| = |f'|$. \\
Let $f'$ be a maximum $(S,T)$-flow of $G(u)$. We define a maximum $(S,T)$-flow $f$
of $G$ as follows:
\begin{equation}
\label{eq:eq2}
f(v,w) =  
  \begin{cases}
      \sum_{x \in in(u)} f'(x,w), & \text{if } v = u \\
      \sum_{x \in out(u)} f'(v,x), & \text{if } w = u \\
      f'(v,w), & \text{otherwise}
   \end{cases} 
\end{equation}
The amount of flow from each $v \in in(u)$ to each $w \in out(u)$ of flow function
$f'$ is redirected over $u$ in $f$. Therefore, $f$ is an valid flow function.
Since $u \notin S\cup T$, it follows that $|f| = |f'|$.
\end{proof}

\begin{figure}
\centering
\includegraphics[width=0.8\textwidth]{../img/network_transformation/clique_expansion.eps}
\caption{Illustration of Lemma \ref{lemma:node_removal} and Equation \ref{eq:eq1} and \ref{eq:eq2}.}
\label{img:clique_expansion}
\end{figure}

In \ShortT{L} all incoming and outgoing edges of a hypernode $v$ have a
capacity equal to $\infty$. The incoming edges are all $\outgoing{e} \in I(v)$ and
the outgoing edges are all $\incoming{e} \in I(v)$. Therefore, we can construct the
following network with Lemma \ref{lemma:node_removal}:

\begin{definition}
Let $T_H$ be a transformation that converts a hypergraph \HypergraphDef~into 
a flow network \ExtendedT{H}{V'} with $V' \subseteq V$. \ShortExT{H}{V'} is defined as follows:
\begin{enumerate}
\item $V_H = V\setminus V' \bigcup\limits_{e \in E}\ \{\incoming{e}, \outgoing{e}\}$
\item $\forall v \in V'$ we add a directed edge $(\outgoing{e_1}, \incoming{e_2}),\ \forall e_1, e_2 \in I(v)$ 
      with $e_1 \neq e_2$ with capacity $\capa_H(\outgoing{e_1}, \incoming{e_2}) = \infty$ (Lemma \ref{lemma:node_removal}).
\item $\forall e \in E$ we add a directed edge $(\incoming{e},\outgoing{e})$
      with capacity $\capa_H(\incoming{e},\outgoing{e}) = \omega(e)$ (same as in \ShortT{L}).
\item $\forall v \in V\setminus V'$ we add for each incident hyperedge $e \in I(v)$ two directed
      edges $(v,\incoming{e})$ and $(\outgoing{e},v)$ with capacity 
      $\capa_H(v,\incoming{e}) = \capa_H(\outgoing{e},v) := \infty$ (same as in \ShortT{L}).
\end{enumerate} 
\end{definition}

An example of the transformation is shown in \autoref{img:heuer_network}.
We have to proof that a minimum-capacity $(S,T)$-cutset
of \ShortExT{H}{V'} is equal with a minimum-weight $(S,T)$-cutset of $H$. However,
we need a preparing lemma in the correctness proof.

\begin{figure}
\centering
\includegraphics[width=0.8\textwidth]{../img/network_transformation/heuer_transformation.eps}
\caption{Transformation of a hypergraph $H$ into an equivalent flow network \ShortExT{H}{V} by removing
         all hypernodes of \ShortT{L}. Note, capacity of the black edges in the flow network is $\infty$.}
\label{img:heuer_network}
\end{figure}

\begin{lemma}
\label{lemma:source_and_sink_removal}
Let $G = (V,E,\capa)$ be a flow network and $f$ a maximum $(S,T)$-flow of $G$.
Further, let $s \in S$ be a source node with $\forall v \in out(s): \capacity{s,v} = \infty$
and $t \in T$ be a sink node with $\forall v \in in(t): \capacity{v,t} = \infty$.
$f_s$ is a maximum $(S',T)$-flow of $G(s)$ and $f_t$ is a maximum $(S,T')$-flow
of $G(t)$ with $S' = S\setminus \{s\} \cup out(s)$ and $T' = T \setminus \{t\}
\cup in(t)$.
\[|f| < \infty \Rightarrow |f| = |f_s| = |f_t|\]
\end{lemma}

\begin{proof}
In Section \ref{sec:applications} we have described how to solve a \emph{multi-source 
multi sink} flow problem by adding a super source node $a$ and super sink node $b$ to the network
and connect $a$ with all sources $s' \in S$ and all sinks $t' \in T$ with $b$.
For $s' \in S: \capacity{a,s'} = \infty$ and for $t' \in T: \capacity{t',b} = \infty$. With Lemma
\ref{lemma:node_removal} follows, that we can remove $s$ from $G$ and insert
a directed edge from $a$ to each $v \in out(s)$ (equal to $G(s)$) and $|f| = |f_s|$. The new flow problem
corresponds to the \emph{multi-source multi-sink} problem with $S'$ and $T$ as source
and sink set. The proof for $G(t)$ is equivalent.
\end{proof}

As a consequence of this lemma, we could replace (or even remove) 
e.g. a source hypernode $v \in S$ of \ShortT{L} and instead add all
incoming hyperedge nodes $\incoming{e} \in I(v)$ as source nodes to the flow 
problem. Because for all incoming resp. outgoing edges of vertices $v$ of 
\ShortT{L} the capacity is $\infty$.

\begin{theorem}
\label{theorem:st_cutset_equal}
A minimum-weight $(S,T)$-cutset of a hypergraph \HypergraphDef~(with $S,T \subseteq V,
S \cap T = \emptyset$) is equivalent with a minimum-capacity $(S',T')$-cutset of the
flow network \ShortExT{H}{V'} ($V' \subseteq V$) with $S' = S \setminus V'\ \cup \bigcup\limits_{e \in N(V' \cap S)} \{\incoming{e}\}$ and 
$T' = T \setminus V'\ \cup \bigcup\limits_{e \in N(V'\cap T)} \{\outgoing{e}\}$.
\label{theorem:heuer_network}
\end{theorem}

\begin{proof}

Applying Lemma \ref{lemma:node_removal} and \ref{lemma:source_and_sink_removal}
on all nodes $v \in V'$ of flow network \ShortT{L} yields to network
\ShortExT{H}{V'} with $S'$ and $T'$ as source and sink sets. A maximum
$(S,T)$-flow $f_L$ of \ShortT{L} is then equal with a maximum $(S',T')$-flow $f_H$
of \ShortExT{H}{V'}. Since $|f_L| < \infty$, only edges between hyperedge 
nodes are contained in a minimum $(S,T)$-cutset of \ShortT{L}. Since $|f_L| = |f_H|$, the same
holds for a minimum $(S',T')$-cutset $E_{min}$ of \ShortExT{H}{V'}. Therefore, $E_{min}$
is equal with a minimum-weight $(S,T)$-cutset of $H$.

%Consider the bipartite graph representation $G_* = (V_*,E_*,c_*)$ 
%of a hypergraph $H = (V,E,c,\omega)$ presented in Section \ref{sec:hypergraph} and \ref{sec:related_lawler}, 
%where for all $v \in V: c_*(v) = \infty$ and for all $e \in E: c_*(e) = 
%\omega(e)$. A minimum-weight $(S,T)$-vertex separator in $G_*$ is equal
%with a minimum-weight $(S,T)$-cutset in $H$. A minimum-weight $(S,T)$-vertex separator can be calculated
%by finding a minimum-capacity $(S,T)$-cutset in \ShortT{L}. Let $G_H$ be the graph obtained by removing
%all $v \in V'\setminus (S \cup T)$ of $G_*$ and adding a clique between all $e \in I(v)$. 
%A minimum-weight $(S,T)$-vertex separator in $G_H$ can be calculated by finding a 
%minimum-capacity $(S,T)$-cutset in our new network \ShortExT{H}{V'\setminus (S \cup T)}.
%We will show that each vertex separator in $G_*$ is also a vertex separator in $G_H$ and
%vice versa. With Lemma \ref{lemma:lemma1} we can show that each minimum $(S,T)$-cutset
%of \ShortExT{H}{V'\setminus (S \cup T)} is equal with a minimum $(S',T')$-cutset of
%\ShortExT{H}{V'} and conclude the proof. We will denote a vertex separator of a graph $G$ with 
%$\V{G}$ and define $V'' := V' \setminus (S \cup T)$. 
%We will show, that $\V{G_*} = \V{G_H}$ with the restriction $\V{G_*} \subseteq E$ and
%$\V{G_H} \subseteq E$.
%
%Assume that $\V{G_*} \subseteq E$ is not a vertex separator in $G_H$. After removing all $e \in \V{G_*}$ of
%$G_H$, there exists still a path $P_H = \{v_1, \ldots, v_k\}$ with $v_1 \in S$ and
%$v_k \in T$ of $G_H$. We can extend $P_H$ to a path $P_*$ in $G_*$.
%We define $P_* := P_H$ and replaces every occurence of a sequence $v_i = e_1 \in E$ and
%$v_{i+1} = e_2 \in E$ with a triple $(e_1,v,e_2)$ in $P_*$, where $v \in e_1 \cap e_2 \cap V''$
%(not empty per construction). $P_*$ does not contain a vertex of $\V{G_*}$, because
%we removed all hyperedge nodes $e \in \V{G_*}$ from $G_H$ before construction of $P_*$ 
%and a hypernode is not part of  the vertex separator $\V{G_*} \subseteq E$ per defintion. 
%$P_*$ connects $S$ and $T$ in $G_*$, which is a contradiction that $\V{G_*}$ 
%is a vertex separator in $G_*$.
%
%Assume that $\V{G_H} \subseteq E$ is not a vertex separator in $G_*$. After removing all $e \in \V{G_H}$ of
%$G_*$, there exists still a path $P_* = \{v_1, \ldots, v_k\}$ with $v_1 \in S$ and
%$v_k \in T$ of $G_*$. We can extend $P_*$ to a path $P_H$ in $G_H$.
%We define $P_H := P_*$ and remove all $v \in P_* \cap V''$ from $P_H$. $G_*$ is a bipartite
%graph per definition. Therefore, each path $P_*$ in $G_*$ is an alternating path of hypernodes and
%hyperedges. The predecessor and successor of a hypernode $v \in P_* \cap V''$ must be hyperedges
%$e_1$ and $e_2$. If $v \in V''$, then $v$ is not contained $G_H$. Instead, there is
%a clique between all $e \in I(v) \Rightarrow$ $(e_1,e_2)$ is contained in $G_H$.
%$P_H$ not contain any vertex of $\V{G_H}$, because we removed all hyperedge nodes $e \in \V{G_H}$ 
%from $G_*$. $P_H$ connects $S$ and $T$ in $G_H$, which is a contradiction that 
%$\V{G_H}$ is a vertex separator in $G_H$.
%
%A minimum-weight $(S,T)$-vertex separator in $G_*$ and $G_H$ contains only hyperedges, because
%the weight of all hypernodes in $G_*$ and $G_H$ is $\infty$. Therefore, each minimum-weight
%$(S,T)$-vertex separator in $G_*$ is also a minimum-weight $(S,T)$-vertex separator in $G_H$,
%because $c(\V{G_*}) = c(\V{G_H})$. With Lemma \ref{lemma:lemma1} follows that
%a minimum-weight $(S,T)$-vertex separator in $G_*$ resp. $G_H$  can also be 
%calculated by finding a minimum-capacity $(S',T')$-cutset in \ShortT{L} resp. \ShortExT{H}{V'}. 
%Therefore, there exists a equivalence between a minimum-weight $(S,T)$-cutset 
%$E_{min}$ of $H$ and the following statements: 
%
%$E_{min}$ is a minimum$\ldots$
%\begin{enumerate}
%\item $\ldots$-weight $(S,T)$-cutset in $H$
%\item $\ldots$-weight $(S,T)$-vertex separator in $G_*$
%\item $\ldots$-capacity $(S,T)$-cutset in \ShortT{L}
%\item $\ldots$-capacity $(S',T')$-cutset in \ShortT{L} (follows from (iii) with Lemma \ref{lemma:lemma1})
%\item $\ldots$-weight $(S,T)$-vertex separator in $G_H$
%\item $\ldots$-capacity $(S,T)$-cutset in \ShortExT{H}{V''}
%\item $\ldots$-capacity $(S',T')$-cutset in \ShortExT{H}{V'} (follows from (vi) with Lemma \ref{lemma:lemma1})
%\end{enumerate}

\end{proof}

Consequently, we can find a minimum-weight $(S,T)$-cutset of $H$ by calculating
a minimum-capacity $(S',T')$-cutset of \ShortExT{H}{V'}. An open problem is how to obtain the 
corresponding minimum-weight $(S,T)$-bipartition. In \ShortT{L} all hypernodes
reachable from source nodes in the residual graph are part of the first and 
all not reachable are part of the second block of the bipartition. Since we removed 
all hypernodes $v \in V'$ in our new network, we have to reconstruct the bipartition
with the following lemma.

\begin{lemma}
\label{lemma:bipartition_construction}
Let $f$ be a maximum $(S,T)$-flow  of \ShortT{L} and $A$ be the set of all nodes reachable
from a node $s \in S$ in the residual graph.
\[ \text{If } v \in A \setminus S \Leftrightarrow \exists e \in I(v): \outgoing{e} \in A \]
\end{lemma}

\begin{proof}
If $\outgoing{e} \in A$, then $v \in A$, because $\capa_L(\outgoing{e},v) = \infty = r_f(\outgoing{e},v)$.
Assume, if $v \in A$, then $\forall e \in I(v): \outgoing{e} \notin A \Rightarrow f(\outgoing{e},v) = 0$. 
Otherwise $r_f(v, \outgoing{e})$ would be greater than zero and this would imply that $\outgoing{e} \in A$
because $v \in A$. Consequently, there is no flow entering node $v$ and due to the conservation
of flow constraint there cannot be any flow leaving node $v$. Therefore, there is no path
from any $s \in S$ to $v$ over a node $\incoming{e}$ because $\forall e \in I(v): \capa_L(\incoming{e},v) = 0 = r_f(\incoming{e},v)$ and no path over $\outgoing{e}$ because $\forall e \in I(v): \outgoing{e} \notin A$
$\Rightarrow$ $v$ is not reachable from any $s \in S$ $\Rightarrow v \notin A$ which is a contradiction
to the assumption that $v \in A$.
\end{proof}

Lemma \ref{lemma:bipartition_construction} gives us an alternative construction for the minimum-weight $(S,T)$-bipartition
of $H$ for both networks \ShortT{L} and \ShortExT{H}{V'}. Regardless of the flow network, we can 
calculate a maximum flow on it and define the set $E''$, which contains all \emph{outgoing hyperedge
nodes} $\outgoing{e}$ \emph{reachable} from a source node $s \in S$ in the \emph{residual graph} of the flow network. 
Further, $(A := \bigcup_{e \in E''} e,\ V\setminus A)$ is a 
minimum-weight $(S,T)$-bipartition of $H$.


\subsection{Low-Degree Hypernodes}
\label{sec:degree_network}

The resulting flow network \ShortExT{H}{V} proposed in Section \ref{sec:heuer_network} has significantly
fewer nodes than the network \ShortT{L} suggested by Lawler. On the other hand, the number of
edges could be much larger. \\
Let's consider a hypernode $v \in V$. We replace $v$ in \ShortT{L} with a clique between all
hyperedges of $I(v)$. The number of edges added to \ShortExT{H}{V} depends on the degree of
$v$. Every hypernode $v \in V$ induce $d(v)(d(v) - 1)$ edges in \ShortExT{H}{V}. In \ShortT{L} a hypernode adds $2d(v)$ edges to the network with the drawback
of an additional node. A simple observation is that for all hypernodes with $d(v) \le 3$ the inequality
$d(v)(d(v) - 1) \le 2d(v)$ holds. Removing such low degree hypernodes not only reduce
the number of nodes, but also the number of edges. \\
Let $V_{d}(n) = \{v \in V\ |\ d(v) \le n\}$ be the set of all hypernodes
with degree smaller or equal $n$. Then our suggested flow network is \ShortExT{H}{V_d(3)}.

\subsection{Removing Graph Hyperedges}
\label{sec:edge_size_network}

If we want to find a minimum-weight $(S,T)$-cutset in a graph $G = (V,E,\omega)$, we do not have to transform
$G$ into an equivalent flow network. We can directly operate on the graph with capacities
$\capacity{e} = \omega(e)$ for all $e \in E$ \cite{ford1956maximal}. Hypergraphs are generalizations of a graph, where
an edge can consist of more than two nodes. However, a hyperedge $e$ of size $2$ can still be 
interpreted as a graph edge. Instead of modeling those edges as described by Lawler \cite{lawler1973}
(see hyperedge $e_2$ in \autoref{img:lawler_transformation}), we can remove all $\incoming{e},\outgoing{e}$ for all $e \in E$
with $|e| = 2$ and add an undirected flow edge between $v_1,v_2 \in e$ (with $v_1 \neq v_2$) with
capacity $\capacity{\{v_1,v_2\}} = \omega(e)$.

\begin{lemma}
\label{lemma:undirected_transformation}
Let $G = (V,E,\capa)$ be an undirected flow network with capacity function $\capa: E \rightarrow \mathbb{N}_+$.
$G$ can be transformed into a directed graph $G'$ such that each maximum $(s,t)$-flow $f$ of $G$
is equal with a maximum $(s,t)$-flow $f'$ of $G'$.
\end{lemma}

\begin{proof}
Assume $\forall e \in E: \capacity{e} = 1$. A maximum $(s,t)$-flow is then equal with the maximum
number of edge-disjoint paths between $s$ and $t$ in a directed graph (Menger's Theorem \cite{menger1927allgemeinen}).
This theorem can also be prooven for undirected graphs if we replace each undirected edge
$e = \{u,v\}$ by five directed edges $(v,\incoming{x}),(w,\incoming{x}),(\incoming{x},\outgoing{x}),(\outgoing{x},v),(\outgoing{x},w)$ (see \autoref{img:menger_transformation})
\cite{menger1927allgemeinen}. Obviously, we can map each set of edge-disjoint paths from $s$ to $t$ from $G'$ to $G$
and vice versa. Therefore, the maximum number of edge-disjoint paths from $s$ to $t$ in
$G'$ is then equal to $G$ and therefore, $|f| = |f'|$. \\
Consider the general case where $\forall e \in E: \capacity{e} \in \mathbb{N}_+$. We can transform the
weighted undirected graph $G$ into an unweighted directed multigraph by replacing each undirected
edge $e = \{u,v\}$ with $\capacity{e}$ undirected edges of weight $1$ (see \autoref{img:menger_transformation}).
Afterwards, we can use the transformation to a unweighted directed multigraph the same way as before.
Again, we can apply Menger's Theorem to show that $|f| = |f'|$. Newman \cite{newman2004analysis}
showed that there is an one-to-one correspondence between a maximum $(s,t)$-flow of an unweighted multigraph
and its corresponding weighted graph where the weight of each edge $(u,v)$ is the number of parallel
edges between $u$ and $v$ of the multigraph.
\end{proof}

\begin{figure}
\centering
\includegraphics[width=0.8\textwidth]{../img/network_transformation/menger_transformation.eps}
\caption{Illustration of the transformation of an unweighted or weighted undirected graph into
         an unweighted or weighted directed graph. The equivalence of a maximum $(s,t)$-flow
         of a unweighted multigraph and their correspondnig weighted graph is a result of 
         a work by Newman \cite{newman2004analysis}.}
\label{img:menger_transformation}
\end{figure}

Note, that each structure of a weighted directed graph illustrated on the right side of
\autoref{img:menger_transformation}~could be transformed into an undirected edge with weight
$\capacity{\{u,v\}} = \capacity{x',x''}$ as a consequence of the construction of the proof of Lemma \ref{lemma:undirected_transformation}.
Each hyperedge $e$ with $|e| = 2$ has exactly this structure in \ShortT{L}. Therefore, we can 
construct the following network:

\begin{definition}
Let $T_G$ be a transformation that converts a hypergraph \HypergraphDef~into 
a flow network \T{G}. \ShortT{G} is defined as follows:
\begin{enumerate}
\item $V_G = V \cup \bigcup\limits_{\substack{e \in E \\ |e| \neq 2}}\ \{\incoming{e}, \outgoing{e}\}$
\item $\forall e \in E$ with $|e| = 2$ and $v_1,v_2 \in e$ ($v_1 \neq v_2$) we add 
      two directed edges $(v_1,v_2)$ and $(v_2,v_1)$ to $E_G$ with capacity $\capa_G(v_1,v_2) = \omega(e)$
      and $\capa_G(v_2,v_1) = \omega(e)$
\item Let $H' = (V,E',c,\omega)$ be the hypergraph with $E' = \{e\ |\ e \in E \land |e| \neq 2\}$,
      then we add all edges of $T_L(H')$ to $E_G$ with their corresponding capacities.
\end{enumerate} 
\end{definition}

An example of transformation \ShortT{G} is shown in \autoref{img:graph_transformation}. A hyperedge
$e$ of size $2$ consists exactly of $4$ nodes and $5$ edges in \ShortT{L} (see \autoref{img:lawler_transformation}).
The same hyperedge induces $2$ nodes and $2$ edges in \ShortT{G}. 
 
\begin{figure}
\centering
\includegraphics[width=0.8\textwidth]{../img/network_transformation/graph_transformation.eps}
\caption{Transformation of a hypergraph into an equivalent flow network by inserting 
         an undirected edge with capacity $\omega(e)$ for each hyperedge of size $2$. 
         Note, capacity of the black edges in the flow network is $\infty$.}
\label{img:graph_transformation}
\end{figure}

\begin{theorem}
A minimum-weight $(S,T)$-cutset of a hypergraph \HypergraphDef~(with $S,T \subseteq V,
S \cap T = \emptyset$) is equal with a minimum-capacity $(S,T)$-cutset of the
flow network \T{G}.
\label{theorem:graph_transformation}
\end{theorem}

\begin{proof}

Before we can apply Lemma \ref{lemma:undirected_transformation} on all hyperedges $e$
with $|e| = 2$, we have to show how to handle the infinite capacity edges of \ShortT{L}.
The flow leaving $\incoming{e}$ is restricted by $\capacity{\incoming{e},\outgoing{e}} = \omega(e)$. Therefore, the flow entering
$\incoming{e}$ is restricted by $f(u,\incoming{e}) + f(v,\incoming{e}) \le \capacity{\incoming{e},\outgoing{e}} = \omega(e)$ with $u \in e$ and
$v \in e$ and $u \neq v \Rightarrow$ $f(u,\incoming{e}) \le \omega(e)$ and $f(v,\incoming{e}) \le \omega(e)$.
The same holds for $f(\outgoing{e},u)$ and $f(\outgoing{e},v)$. Therefore, we can replace each infinite 
capacity of an edge entering $\incoming{e}$ or leaving $\outgoing{e}$ with $\omega(e)$ without changing the amount
of a maximum $(S,T)$-flow. We call the capacity adapted network \ShortT{L'}.\\
Applying the transformation of Lemma \ref{lemma:undirected_transformation} on each undirected
edge of \ShortT{G} results in flow network \ShortT{L'}. It follows, that a maximum $(S,T)$-flow of \ShortT{G} is equal
with a maximum $(S,T)$-flow of \ShortT{L'} and \ShortT{L}. Consequently, a minimum-capacity
$(S,T)$-cutset of \ShortT{G} is equal with a minimum-weight $(S,T)$-cutset of $H$.

%We define the bijective function $\Phi: E_L \rightarrow E_G$ as follows
%\[ \Phi(\incoming{e},\outgoing{e}) = 
%   \begin{cases}
%      (\incoming{e},\outgoing{e}), \text{if } |e| \neq 2, \\
%      \{v_1,v_2\}, \text{otherwise (with $v_1,v_2 \in e$ and $v_1 \neq v_2$)}
%   \end{cases} \]
%
%We will show that each $(S,T)$-cutset $A_L$ of \ShortT{L} is a $(S,T)$-cutset $\Phi(A_L)$ of
%\ShortT{G} and vice versa. Per defintion $\capa_L(A_L) = \capa_G(\Phi(A_L))$ and for each $(S,T)$-cutset
%$A_G$ of \ShortT{G} $\capa_G(A_G) = c_L(\Phi^{-1}(A_G))$. Therefore, each minimum-capacity $(S,T)$-cutset
%of \ShortT{L} must be a minimum-capacity $(S,T)$-cutset of \ShortT{G} and vice versa. In the following 
%let $E^* = \bigcup_{e \in E} \{(\incoming{e},\outgoing{e})\}$. \\
%Let $A_L \subseteq E^*$ be a $(S,T)$-cutset of \ShortT{L}. Assume $\Phi(A_L)$ is not a $(S,T)$-cutset
%of \ShortT{G}. Then there exists a path
%$P_G = \{v_1,\ldots,v_k\}$ connecting $S$ and $T$ in \ShortT{G} not containing any edge $e \in \Phi(A_L)$.
%Let $P_L$ be the path in \ShortT{L} obtained by adding edge $\Phi^{-1}(v_i,v_{i+1}) = (\incoming{e},\outgoing{e})$ between all
%$v_i \in V$ and $v_{i+1} \in V$ in $P_G$. $\Phi^{-1}(v_i,v_{i+1}) \notin A_L$, because
%$P_G$ does not contain any edge of $\Phi(A_L)$. $P_L$ connects $S$ and
%$T$ in \ShortT{L}, which is a contradiction to the assumption that $A_L$ is a $(S,T)$-cutset. \\
%Let $A_G \subseteq \Phi(E^*)$ be a $(S,T)$-cutset in \ShortT{G}. Let's assume $\Phi^{-1}(A_G)$ is not
%a $(S,T)$-cutset of \ShortT{L}. Then there exists a path $P_L = \{v_1,\ldots,v_k\}$ 
%connecting $S$ and $T$ in \ShortT{L} not containing any
%edge $e \in \Phi^{-1}(A_G)$. Let $P_G$ be the path in \ShortT{G} obtained by removing each edge
%$(v_i,v_{i+1})$ with $v_i = \incoming{e}$ and $v_{i+1} = \outgoing{e}$ and $|e| = 2$ from $P_L$. Based on the construction
%of \ShortT{L} the predecessor of $v_i$ and successor of $v_{i+1}$ must be hypernodes $v_1,v_2 \in e$.
%Therefore, $P_G$ is a valid path in \ShortT{G} connecting $S$ and $T$ and not contain any edge of 
%$A_G$. This is a contradiction to the assumption that $A_G$ is a $(S,T)$-cutset.
\end{proof} 

A minimum-weight $(S,T)$-cutset of $H$ could also be calculated with \ShortT{G}. 
Each edge $(v_1,v_2)$ with $v_1,v_2 \in V$ of the minimum-capacity
$(S,T)$-cutset of \ShortT{G} can be mapped to their corresponding hyperedge. 
Since there exists a one-one correspondence between the hypernodes
of \ShortT{L} and \ShortT{G} the corresponding bipartition are all hypernodes \emph{reachable}
from all nodes in $S$ and all not \emph{reachable} from $S$ in the \emph{residual graph}
of \ShortT{G}. 

\subsection{Combining Techniques}
\label{sec:hybrid_network}

In many real-world instances, the average hyperedge size and hypernode degree are inversely
proportional to each other. E.g., if the number of hyperedges is significantly larger than the
number of hypernodes the average hypernode degree is usually much larger than $3$. Whereas
the average hyperedge size is often equal to $2$. If the number of hyperedges is nearly equal
to the number of hypernodes, the average hypernode degree is usually smaller or equal than $3$. Whereas
the average hyperedge size is often much larger than $2$. Of course, we can construct instances
where this inversely proportional relationship cannot be observed, but in many real-world instances,
we often find the described behavior. \\
Currently, we have two different modeling approaches which either perform better on low degree
hypernode instances or small hyperedge size instances. Taking our observation from real-world instances 
into account means that either \ShortT{G} or \ShortExT{H}{V_d(3)} performs significantly better on
a specific instance. It would be preferable to combine the two approaches into one network 
which performs on most instances best. \\

\begin{definition}
Let $T_{\text{Hybrid}}$ be a transformation that converts a hypergraph \HypergraphDef~into 
a flow network \Hybrid, where $V' = \{v \in V_d(3)\ |\ \forall e \in I(v): |e| \neq 2\}$. 
\ShortHybrid~is defined as follows:
\begin{enumerate}
\item $V_{\text{Hybrid}} = V\setminus V' \bigcup\limits_{\substack{e \in E \\ |e| \neq 2}}\ \{\incoming{e}, \outgoing{e}\}$
\item $\forall v \in V'$ we add a directed edge $(\outgoing{e_1}, \incoming{e_2}),\ \forall e_1, e_2 \in I(v)$ 
      ($e_1 \neq e_2$) with capacity $\capa_{\text{Hybrid}}(\outgoing{e_1}, \incoming{e_2}) = \infty$ (Lemma \ref{lemma:node_removal}).
\item $\forall e \in E$ with $|e| = 2$ and $v_1,v_2 \in e$ ($v_1 \neq v_2$) we add 
      two directed edges $(v_1,v_2)$ and $(v_2,v_1)$ with capacity $\capa_{\text{Hybrid}}(v_1,v_2) = \omega(e)$
      and $\capa_{\text{Hybrid}}(v_2,v_1) = \omega(e)$ (Lemma \ref{lemma:undirected_transformation})
\item $\forall e \in E$ with $|e| \neq 2$ we add a directed edge $(\incoming{e},\outgoing{e})$
      with capacity $\capa_{\text{Hybrid}}(\incoming{e},\outgoing{e}) = \omega(e)$ (same as in \ShortT{L}).
\item $\forall v \in V\setminus V'$ we add for each incident hyperedge $e \in I(v)$ with $|e| \neq 2$ 
      two directed edges $(v,\incoming{e})$ and $(\outgoing{e},v)$ with capacity 
      $\capa_{\text{Hybrid}}(v,\incoming{e}) = \capa_{\text{Hybrid}}(\outgoing{e},v) := \infty$ (same as in \ShortT{L}).
\end{enumerate} 
\end{definition}

\autoref{img:transformation_chain} summarizes all explained transformations of this section.
We can prove the correctness of \ShortHybrid~with Lemma \ref{lemma:node_removal}, \ref{lemma:source_and_sink_removal}
and \ref{lemma:undirected_transformation} as used in the proof of Theorem \ref{theorem:heuer_network}
and \ref{theorem:graph_transformation}. A minimum-weight $(S,T)$-cutset of $H$ is equal with a
minimum-capacity $(S',T')$-cutset of \ShortHybrid. \\
Per definition of \ShortHybrid~we prefer a hyperedge removal over a hypernode removal. E.g., if
a hypernode has a degree smaller or equal than $3$, we only remove it, if there is no hyperedge
$e \in I(v)$ with $|e| = 2$. The reason is that a hyperedge removal always decreases the number of nodes
and edges more than a hypernode removal. \\
\begin{figure}[ht!]
\centering
\includegraphics[width=0.95\textwidth]{../img/network_transformation/hybrid_transformation.eps}
\caption{Illustration of all presented sparsifying techniques of the flow network of a hypergraph.}
\label{img:transformation_chain}
\end{figure}
The minimum-weight $(S,T)$-cutset of $H$ can be calculated with the same technique described in Section
\ref{sec:edge_size_network}. Let $(A,V\setminus A)$ be the corresponding bipartition.
$A$ is the union of all reachable hypernodes from $S'$ and the union of
all reachable \emph{outgoing hyperedge nodes} $\outgoing{e}$ from $S'$ (see Section \ref{sec:heuer_network} 
and Lemma \ref{lemma:bipartition_construction}). 


\newpage

%%%%%%%%%%%%%%%%%%%%%%%%%%%%%%%%%%%%%%%%%%%%%%%%%%%%%%%%%%%%%%%%%%%%%%

\section{Using Max-Flow-Min-Cut Computations as a Local Search Strategy}

We will give now a detailed description of our flow-based refinement framework. The main
idea is to extract a subhypergraph $H_{V'}$ out of the original hypergraph $H$, which is already
partitioned into $k$ blocks. $V'$ is choosen in such a way that it is a subset of two
adjacent blocks $V_i$ and $V_j$. We will show how to configure
the sources $S$ and sinks $T$ of the corresponding flow network such that
a minimum $(S,T)$-bipartition of $H_{V'}$ improves the connectivity metric on $H$
(see Section \ref{sec:source_and_sink}). Further, we describe how the ideas of
Sanders and Schulz \cite{sanders2011engineering} (see Section \ref{sec:flow_local_search_graph}) 
could be adapted to work in a $n$-level hypergraph partitioner, called \emph{KaHyPar}
(see Section \ref{sec:mbmc_hypergraphs} and \ref{sec:flow_local_search_hypergraph}). 

\subsection{Modelling Sources and Sinks}
\label{sec:source_and_sink}

Let $H = (V,E,c,\omega)$ be a hypergraph and $B_1 := (V_1, V_2)$ be a bipartition.
$H_{V'} = (V',E_{V'},c,\omega)$ is the subhypergraph
induced by $V' \subseteq V$ with $E_{V'} = \{ e \cap V'\ |\ e \in E: e \cap V' \neq \emptyset\}$. 
Further, let $E_{\emptyset} = \{ e \cap V'\ |\ e \in E: e \cap V' = \emptyset\} $
be the set of all hyperedges contained in $H$, but not in $H_{V'}$. $T_L(H_{V'})$ 
(see Section \ref{sec:related_lawler}) is the flow problem induced by $H_{V'}$ with a
source set $S$ and a sink set $T$. Let $(V_1',V_2')$ be the minimum $(S,T)$-bipartition
obtained by a maximum $(S,T)$-flow calculation on $T_L(H_{V'})$ with $f$ as maximum flow
function. We can extend the bipartition $(V_1',V_2')$ of $H_{V'}$ to a bipartition 
$B_2 := (V_1 \setminus V'\ \cup\ V_1',\ V_2 \setminus V'\ \cup V_2' )$ on $H$. Finally,
we need to define the cut on hypergraph $H$ and its subhypergraph $H_{V'}$ related to a
bipartion $(V_1,V_2)$. For this porpose we define the set of all cut hyperedges
$E(V_1,V_2) := \{e \in E\ |\ \exists u,v \in e: u \in V_1\ \land\ v \in V_2\}$.
We define the cut on $H$ resp. $H_{V'}$ with
\begin{align*}
\omega_H(V_1,V_2) := \sum_{e \in E(V_1,V_2)} \omega(e) \\
\omega_{H_{V'}}(V_1,V_2) := \sum_{e \in E(V_1,V_2) \setminus E_{\emptyset}} \omega(e) 
\end{align*}
Some will be wondering about the defnition of the cut $\omega_{H_{V'}}$ of $H_{V'}$ over the
cut edges of $H$. A cut hyperedge $e$ of $H$ must not necessarily be a cut hyperedge
of $H_{V'}$. E.g., if $e = \{v_1,v_2\}$ with $v_1 \in V_1$ and $v_2 \in V_2$, but
$v_1 \in V'$ and $v_2 \notin V'$. Then $e$ is cut in $H$, but not in $H_{V'}$, because
$v_2$ is removed from $e$ per definition of $E_{V'}$. However, the reason that we still
define $e$ as cut hyperedge of $H_{V'}$ has to do with our problem statement, 
which we will define as follows:

\begin{problem}
\label{prob:ST}
How do we have to define the source set $S$ and sink set $T$ for a subhypergraph $H_{V'}$ 
(with $V' \subseteq V$) and a bipartition $B_1$, such that 
after a maximum $(S,T)$-flow calculation (with $f$ as maximum flow function)
the resulting bipartition $B_2$ on $H$ satisfy the following conditions:
\begin{enumerate}
\item $\omega_H(B_2) \le \omega_H(B_1)$
\item $\Delta_{H} := \omega_H(B_1) - \omega_H(B_2) = \omega_{H_{V'}}(B_1) - |f| =: \Delta_{H_{V'}}$
\end{enumerate}
\end{problem}

The first condition ensures that a maximum $(S,T)$-flow calculation on $T_L(H_{V'})$ never 
decrease the cut of $H$. The existence of the second condition has practical reasons. First, we
can simply update the cut metric via $\omega_H(B_2) = w_H(B_1) - \Delta_{H_{V'}}$,
instead of summing up the weight of all cut hyperedges. Since, we have to setup the subhypergraph
$H_{V'}$ before each maximum flow computation we can implicitly calculate $\omega_{H_{V'}}(B_1)$.
Therefore, the cut metric can be updated after the \emph{Max-Flow-Min-Cut} computation
in constant time instead of \BigO{|E|}. On the other hand, we can assert the correctness of
our own maximum flow algorithm. If $\Delta_H \neq \Delta_{H_{V'}}$, then with high probability our
flow algorithm is incorrect. Also, the reason why we define $\omega_{H_{V'}}(V_1,V_2)$ over
the cut hyperedges of $H$ is due to the fact that the equality
\[\Delta_{H} := \omega_H(B_1) - \omega_H(B_2) = \omega_{H_{V'}}(B_1) - \omega_{H_{V'}}(B_2)\]
holds. If we are able to show that $|f| = \omega_{H_{V'}}(B_2)$, we simultanously show
that $\Delta_H = \Delta_{H_{V'}}$.\\
We will now present a solution for our problem statement. First, we show how $S$ and $T$
can be choosen to satisfy condition (i). Afterwards, we extend $S$ and $T$ with additional
nodes to fullfil condition (ii). Finally, we show how $S$ and $T$ can be modified, such that
we can obtain smaller cuts on $H$ and simultanously satisfy condition (i) and (ii) of our problem
statement. \\
Let $V' \subseteq V$ and $\delta B = \{ e \in E\ |\ \exists u,v \in e: u \in V'\ \land\ v \notin V' \}$
be the set of all \emph{Border Hyperedges}. For a bipartition $(V_1,V_2)$ of $H$, we say
$v \in V_1$ is a source node of the flow network $T_L(H_{V'})$, if there exists
a hyperedge $e \in \delta B$ containing $v$ and at least one other node $u \in V_1$ with
$u \notin V'$. More formal:
\begin{align}
%S = \bigcup\limits_{\substack{e \in \delta B \\ |e \setminus V' \cap V_1| \neq 0}} (e \cap V_1 \cap V') \\
S_1 = \{ s \in V' \cap V_1\ |\ \exists v \notin V': \exists e \in \delta B: v \in V_1\ \land\ s,v \in e \} \label{S1_border_hyperedges}\\
T_1 = \{ t \in V' \cap V_2\ |\ \exists v \notin V': \exists e \in \delta B: v \in V_2\ \land\ v,t \in e \} \label{T1_border_hyperedges}
\end{align}
An example of a \emph{Max-Flow-Min-Cut} computation on $H_{V'}$ with $S$ and $T$ as source and
sink set is illustrated in \autoref{img:border_hyperedges}.
\begin{figure}
\centering
\includegraphics[width=0.95\textwidth]{../img/source_and_sink_set/border_hyperedges.eps}
\caption{Example how \emph{Border Hyperedges} are modelled as sources and sinks.}
\label{img:border_hyperedges}
\end{figure}

\begin{lemma}
\label{cut_decrease_proof}
Let $B_1$ be a bipartition of $H$ and $T_L(H_{V'})$ the flow network of subhypergraph
$H_{V'}$ with $S$ and $T$ as defined in Equation \ref{S1_border_hyperedges} and \ref{T1_border_hyperedges} (with $V' \subseteq V$).
Let $B_2$ be the bipartition obtained by a maximum $(S,T)$-flow computation on $T_L(H_{V'})$.
Then, $\omega_H(B_2) \le \omega_H(B_1)$.
\end{lemma}

\begin{proof}
A $(S,T)$ \emph{Max-Flow-Min-Cut} computation on $T_L(H_{V'})$ yields to a minimum $(S,T)$-cutset on 
$H_{V'}$ \cite{ford1956maximal}. Thus, for all hyperedges $e \notin \delta B \cup E_{\emptyset}$, which
are cut in $B_2$, the sum of their weight must be less or equal than the sum of all cut hyperedges
$e \notin \delta B \cup E_{\emptyset}$ of bipartition $B_1$. We need to show that a non-cut
hyperedge $e \in \delta B$ of $B_1 = (V_1,V_2)$ cannot become a cut hyperedge of
$B_2 = (V_1',V_2')$. Let $e \in \delta B$ be such a hyperedge. $e$ must be either a subset of $V_1$ or $V_2$, otherwise
$e$ is a cut hyperedge. Let $e \subseteq V_1$, then $e \cap V' \subseteq S$ (see Equation \ref{S1_border_hyperedges}). 
Defining a node $s \in S$ as source node means that it cannot change its block after a \emph{Max-Flow-Min-Cut}
computation. Therefore, $e \subseteq V_1$ and $e \subseteq V_1' \Rightarrow e$ is a non-cut
hyperedge in $B_2$. The proof for $e \subseteq V_2$ is equivalent $\Rightarrow \omega_H(B_2) 
\le \omega_H(B_1)$.
\end{proof} 

In the next step we will show how $S$ and $T$ can be extended to fullfil condition (ii)
of Problem \autoref{prob:ST}. Currently, $|f| \le \omega_{H_{V'}}(B_2)$ (without a prove).
Obviously, some nodes are missing in $S$ and $T$. Do understand which nodes are missing consider
\autoref{img:non_cut_flow_hyperedges}. Transformation $1$ illustrates our current modelling
approach defined in Equation \ref{S1_border_hyperedges} and \ref{T1_border_hyperedges}. The maximum
flow on this network is $|f| = 1$, but the resulting minimum $(S,T)$-bipartition $B_2$ induced 
a cut of $\omega_{H_{V'}}(B_2) = 2$. This due to the fact that $e_1$ and $e_3$ are cut
hyperedges in $H$, but non-cut hyperedges in $H_{V'}$. The actual cut of $H_{V'}$ is therefore
$1$ (instead of $2$) and this is also a minimum $(S,T)$-cut. Transformation $2$ illustrates the
correct modelling approach for cut hyperedges of $H$ which are non-cut hyperedges in $H_{V'}$.
For each hyperedge $e \in \delta B$ with $e \cap V' \subseteq V_2$ and $e \setminus V' \cap V_1 \neq \emptyset$, 
we add the \emph{incomming hyperedge node} $e'$ to $S$.
More formal:
\begin{align}
S = S_1 \cup \{e'\ |\ e \cap V' \subseteq V_2 \ \land\ e \setminus V' \cap V_1  \neq \emptyset\} \label{S_border_hyperedges}\\
T = T_1 \cup \{e''\ |\ e \cap V' \subseteq V_1 \ \land\ e \setminus V' \cap V_2 \neq \emptyset \} \label{T_border_hyperedges}
\end{align}

\begin{figure}
\centering
\includegraphics[width=1.0\textwidth]{../img/source_and_sink_set/non_cut_flow_hyperedges.eps}
\caption{In this example $e_1$ and $e_3$ are cut hyperedges of the hypergraph, but not of
        the subhypergraph induced by the flow problem. Modelling the \emph{outgoing} resp.
        \emph{incomming} hyperedge node of $e_1$ resp. $e_2$ as sink resp. source ensures
        that $\Delta_H = \Delta_{H_{V'}}$.} 
\label{img:non_cut_flow_hyperedges}
\end{figure}

\begin{lemma}
\label{lemma:delta_proof}
Let $B_1$ be a bipartition of $H$ and $T_L(H_{V'})$ the flow network of subhypergraph
$H_{V'}$ with $S$ and $T$ as defined in Equation \ref{S_border_hyperedges} and \ref{T_border_hyperedges} (with $V' \subseteq V$).
Let $B_2$ be the bipartition obtained by a maximum $(S,T)$-flow computation on $T_L(H_{V'})$
with $f$ as maximum flow function.
Then, $\omega_{H_{V'}}(B_2) = |f|$ ($\Rightarrow \Delta_H = \Delta_{H_{V'}}$).
\end{lemma}
 
\begin{figure}
\centering
\includegraphics[width=0.95\textwidth]{../img/source_and_sink_set/delta_proof_illustration.eps}
\caption{Illustration of the two cases presented in proof of Lemma \ref{lemma:delta_proof} in
         order to remove a hypernode $v \in V'' \cap S$ from $T_L(H_{V' \cup V''})$. }
\label{img:delta_proof_illustration}
\end{figure}

\begin{proof}
Let $V'' = \bigcup_{e \in \delta B} e \setminus V'$ be the set of all hypernodes contained
in a \emph{border hyperedge}, but not in $V'$. Let $H_{V' \cup V''}$ be the subhypergraph
obtained by extending $H_{V'}$ with all missing hypernodes of \emph{border hyperedges}. 
We define the flow problem $T_L(H_{V' \cup V''})$ with $S' = S_1\ \cup\ (V'' \cap V_1)$ and 
$T' = T_1\ \cup\ (V'' \cap V_2)$ as sources and sinks. Further, let $f'$ be a maximum $(S',T')$-flow 
on $T_L(H_{V' \cup V''})$ and $B_2$ be the corresponding minimum $(S',T')$-bipartition. Because
all hypernodes which are part of a hyperedge in $H$ and also in $H_{V'}$ are fully contained
in $H_{V' \cup V''}$ the equality $|f'| = \omega_{H_{V'}}(B_2)$ holds. In the following we 
present a technique with which we can obtain a new flow network $T_L(H_{V' \cup V''\setminus \{v\}})$
with $v \in V''$ and $S''$ and $T''$ as sources and sinks. Simultanously we map
the maximum $(S',T')$-flow $f'$ of $T_L(H_{V' \cup V''})$ to a maximum $(S'',T'')$-flow
of $T_L(H_{V' \cup V''\setminus \{v\}})$ with $|f'| = |f''|$. Applying this technique successively
on all nodes $v \in V''$ will result in the flow network $T_L(H_{V'})$ with $S$ and $T$ as sources and sinks
from Equation \ref{S_border_hyperedges} and \ref{T_border_hyperedges}. \\
A hypernode $v \in V''$ is either a source or a sink. We will show how to remove a source
hypernode $v \in V'' \cap S'$. At beginning we define $S'' := S'$, $T'' := T'$ and 
$f'' := f'$. In order to remove $v \in V''$ we have to distingush two cases 
based on a incident hyperedge $e \in I(v)$ of $v$: 

$\mathbf{e \cap S \setminus \{v\} \neq \emptyset}$:
Then there exists a hypernode $u \in e \cap S$ with $u \neq v$. We define
$f''(u,e') = f''(u,e') + f'(v,e')$ and $f''(s,u) = f''(s,u) + f'(v,e')$. 

$\mathbf{e \cap S \setminus \{v\} = \emptyset}$:
In this case $e$ must be a cut hyperedge in $H$, but not in $H_{V'}$, otherwise
there would exist a hypernode $u \in e \cap S$ (see Equation \ref{S1_border_hyperedges}).
We define $S'' = S'' \cup \{e'\}$. Simultanously, we set  
$f''(s,e') = f'(v,e')$. 

The two cases are illustrated in \autoref{img:delta_proof_illustration}.
We can remove $v$ from $T_L(H_{V' \cup V''})$ after applying this procedure for 
all $e \in I(v)$. The cases for a hypernode $v \in V'' \cap T'$ are equivalent. 
$f''$ is a valid flow function per construction and $|f'| = |f''|$. Also $f''$
is maximum $(S'',T'')$-flow on $T_L(H_{V' \cup V''\setminus \{v\}})$, otherwise
we can map a augmenting path in the residual graph $T_L(H_{V' \cup V''\setminus \{v\}})$
to a augmenting path in $T_L(H_{V' \cup V''})$ (without a proof). We can successively
remove all $v \in V''$ from $T_L(H_{V' \cup V''})$ with this method. \\
The resulting  flow network is $T_L(H_{V'})$. For each $e \in E$ which is cut in $H$, but not
in $H_{V'}$, we have added the corresponding \emph{incomming hyperedge node} $e'$
or \emph{outgoing hyperedge node} $e''$ to $S''$ resp. $T''$. Therefore, 
$S''$ and $T''$ are equal with $S$ and $T$ defined in Equation \ref{S_border_hyperedges}
and \ref{T_border_hyperedges}. Finally, the flow function $f''$ is a maximum
$(S,T)$-flow on $T_L(H_{V'})$ and $|f''| = |f'| = \omega_{H_{V'}}(B_2)$ per construction.
\end{proof}

With our current modelling approach we are able to satisfy all conditions of
our problem statement. However, sometimes we define hypernodes as source resp.
sink which are unnecessary. For a explanation consider \autoref{img:cut_border_hyperedges}.
Hyperedge $e_1$ is cut in $H_{V'}$ and contains hypernodes from both
blocks, which are not in the flow problem. Regardless what we do in $H_{V'}$
we can not remove $e_1$ from cut in $H$. Using our suggested source and sink modelling has
as consequence that $e_1$ and $e_2$ are still cut after a \emph{Max-Flow-Min-Cut} 
computation (see \emph{Transformation 1} in \autoref{img:cut_border_hyperedges}). Another approach
is to define for hyperedges which are cut of $H_{V'}$ and are also of $H$ the \emph{incomming}
resp. \emph{outgoing hyperedge node} as source resp. sink (see \emph{Transformation 2} in
\autoref{img:cut_border_hyperedges}). In our example all hypernodes of $e_1$ are still able 
to move and a \emph{Max-Flow-Min-Cut} computation removes $e_2$ from cut.\\
To define our final source and sink set, we split the set of all \emph{border hyperedges} into
three different disjoint subsets as follows:

\begin{enumerate}
\item $\delta B_1 = \{ e \in \delta B \ |\ e \subseteq V_1\ \lor\ e \subseteq V_2 \}$
\item $\delta B_2 = \{ e \in \delta B \ |\ e \cap V' \not\subseteq V_1\ \land\ e \cap V' \not\subseteq V_2\}$
\item $\delta B_3 = \{ e \in \delta B \setminus \delta B_1 \ |\ (e \cap V' \subseteq V_1\ \lor\ e \cap V' \subseteq V_2 \}$
\end{enumerate}

$\delta B_1$ contains all non-cut \emph{border hyperedges} of $H$. $\delta B_2$ contains 
all \emph{cut border hyperedges} of $H$, which are also cut in $H_{V'}$ and $\delta B_3$ contains
all \emph{cut border hyperedges} of $H$, which are non-cut in $H_{V'}$.
\begin{align}
S = \bigcup\limits_{\substack{e \in \delta B_1 \\ e \subseteq V_1}} e\cap V'\ \cup \bigcup\limits_{\substack{e \in \delta B_2 \cup \delta B_3 \\ e \setminus V' \cap V_1 \neq \emptyset}} \{e'\} \label{S_final_border_hyperedges}\\
T = \bigcup\limits_{\substack{e \in \delta B_1 \\ e \subseteq V_2}} e\cap V'\ \cup \bigcup\limits_{\substack{e \in \delta B_2 \cup \delta B_3 \\ e \setminus V' \cap V_2 \neq \emptyset}} \{e''\} \label{T_final_border_hyperedges}
\end{align}

\begin{figure}
\centering
\includegraphics[width=0.95\textwidth]{../img/source_and_sink_set/cut_border_hyperedges.eps}
\caption{Example how \emph{Cut Border Hyperedges} are modelled as sources and sinks. In this
         example $e_1$ contains node from block $V_1$ and $V_2$ not contained in the flow problem. Therefore,
         we can not remove $e_1$ from cut. Treating $e_1$ as a \emph{Border Hyperedge} would result
         in Transformation $1$. This has the consequence that we are not able to remove $e_2$
         from cut with a \emph{Max-Flow-Min-Cut} computation. Defining the \emph{incomming} resp.
         \emph{outgoing} hyperedge of $e_1$ as source resp. sinks allows the corresponding hypernodes
         of $e_1$ still to move. The consequence is that we can remove $e_2$ from cut with a
         \emph{Max-Flow-Min-Cut} computation in Transformation $2$.}
\label{img:cut_border_hyperedges}
\end{figure}

Equation \ref{S_final_border_hyperedges} and \ref{T_final_border_hyperedges} 
are illustrated in \autoref{img:st_modelling_summary}.
A \emph{Max-Flow-Min-Cut} computation on $T_L(H_{V'})$ with $S$ and $T$ as defined
in Equation \ref{S_final_border_hyperedges} and \ref{T_final_border_hyperedges} satisfy
condition (i) and (ii) of Problem \ref{prob:ST}. This can be prooven with similiar techniques
used in the proof of Lemma \ref{cut_decrease_proof} and \ref{lemma:delta_proof}. A maximum
$(S,T)$-flow calculation yields to a minimum $(S,T)$-cut on $H_{V'}$. A non-cut hyperedge
$e \in \delta B_1$ can not become a cut hyperedge after a \emph{Max-Flow-Min-Cut} computation,
because we still define all hypernodes $v \in e \cap V'$ as sources resp. sinks. Therefore,
$\omega_H(B_2) \le \omega_H(B_1)$. We can proof Lemma \ref{lemma:delta_proof} for our
new source and sink set if we adapt the conditions of the cases for a hyperedge $e \in I(v)$ 
based on the set $\delta B_1$, $\delta B_2$ and $\delta B_3$ where $e$ is contained. If
$e \in \delta B_1$, then there must exists a hypernode $u \in e \cap S \setminus \{v\}$ on
which we apply the first case (Case $1$: $e \cap S \setminus \{v\} \neq \emptyset$). 
For all $e \in \delta B_2 \cup \delta B_3$, we simply apply the second case
(Case $2$: $e \cap S \setminus \{v\} = \emptyset$). After removing all hypernodes $v \in V''$
the resulting network is $T_L(H_{V'})$ with $S$ and $T$ as defined in Equation
\ref{S_final_border_hyperedges} and \ref{T_final_border_hyperedges}. Further, the flow function
$f''$ is a maximum $(S,T)$-flow on $T_L(H_{V'})$ with $|f''| = |f'| = \omega_{H_{V'}}(B_2)
\Rightarrow \Delta_H = \Delta_{H_{V'}}$. \\
Finally, we want to show that for a minimum $(S',T')$-bipartition $B_2$ with $S'$ and $T'$ as defined
in Equation \ref{S_final_border_hyperedges} and \ref{T_final_border_hyperedges} and a minimum
$(S,T)$-bipartition $B_3$ with $S$ and $T$ as defined in Equation \ref{S_border_hyperedges}
and \ref{T_border_hyperedges} calculated with flow network $T_L(H_{V'})$ the inequality 
$\omega_H(B_2) \le \omega_H(B_3)$ holds. For this propose we need a preparing lemma.

\begin{lemma}
\label{lemma:ST_subset}
Let $G = (V,E,c)$ be a flow network with sources $S$ and sinks $T$. Further, let $S' \subseteq S$
and $T' \subseteq T$. The value of a maximum $(S',T')$-flow $f'$ is less or equal than the value
of a maximum $(S,T)$-flow $f$. More formal, $|f'| \le |f|$.
\end{lemma}

\begin{proof}
Assume $|f'| > |f|$. Then, we can simply set $f = f'$, because $S' \subseteq S$ and
$T' \subseteq T$. But this is a contradiction to assumption that $f$ is a maximum $(S,T)$-flow
on $G$. Therefore, $|f'| \le |f|$.
\end{proof}

In the following theorem, we denote with $S$ and $T$ the source and sink sets as defined in 
Equation \ref{S_border_hyperedges} and \ref{T_border_hyperedges} and with $S'$ and $T'$ the
source and sink sets as defined in Equation \ref{S_final_border_hyperedges} and
\ref{T_final_border_hyperedges}.

\begin{theorem}
Let $H$ be a hypergraph and $H_{V'}$ be the subhypergraph induced by the subset $V' \subseteq V$.
Further, $B_1$ is the current bipartition of $H$. For a minimum $(S',T')$-bipartition $B_2$ and
a minimum $(S,T)$-bipartition $B_3$ obtained by a maximum $(S',T')$- resp. $(S,T)$-flow calculation
on $T_L(H_{V'})$ the inequality $\omega_H(B_2) \le \omega_H(B_3) \le \omega_H(B_1)$ holds.
\end{theorem}

\begin{proof}
Let $(\bar{S'},\bar{T'})$ resp. $(\bar{S},\bar{T})$ be the sets obtained by removing all
\emph{incomming} and \emph{outgoing hyperedge nodes} $e'$ and $e''$ from $(S',T')$ resp.
$(S,T)$. It holds that $\bar{S'} \subseteq \bar{S}$ and $\bar{T'} \subseteq \bar{T}$. Afterwards,
we extend the subhypergraph $H_{V'}$ with all hypernodes $V'' = \bigcup_{e \in \delta B} e \setminus V'$
and obtain subhypergraph $H_{V' \cup V''}$ with flow network $T_L(H_{V' \cup V''})$.
Also we extend $(\bar{S'},\bar{T'})$ and $(\bar{S},\bar{T})$ exactly in the same way as 
in the proof of Theorem \ref{lemma:delta_proof}. With the \emph{Max-Flow-Min-Cut}-Theorem 
\cite{ford1956maximal} we can conclude that the cut value $\omega_{H_{V'}}(B_2)$ of 
a minimum $(\bar{S'},\bar{T'})$-bipartition $B_2$ on $H_{V'}$ is equal with the value of a 
maximum $(\bar{S'},\bar{T'})$-flow $f'$ on $T_L(H_{V' \cup V''})$. 
The same holds for a minimum $(\bar{S},\bar{T})$-bipartition
$B_3$ and a maximum $(\bar{S},\bar{T})$-flow $f$. After extending $(\bar{S'},\bar{T'})$ 
resp. $(\bar{S},\bar{T})$ with all hypernodes of $V''$ $\bar{S'} \subseteq \bar{S}$
and $\bar{T'} \subseteq \bar{T}$ still holds. With Lemma \ref{lemma:ST_subset} and the 
\emph{Max-Flow-Min-Cut}-Theorem follows $\omega_{H_{V'}}(B_2) = |f'| \le |f| = \omega_{H_{V'}}(B_3)$.\\
We can transform $(\bar{S'},\bar{T'})$ resp. $(\bar{S},\bar{T})$ and flow network $T_L(H_{V'\cup V''})$
back to $T_L(H_{V'})$ with $(S',T')$ resp. $(S,T)$ as source and sink sets
with the technique described in the proof of Theorem \ref{lemma:delta_proof} and in the sketch
of the proof for our new source and sink sets (see Equation \ref{S_final_border_hyperedges} and
\ref{T_final_border_hyperedges}). Therefore, the inequality still holds for bipartitions
$B_2$ and $B_3$ obtained by a maximum $(S',T')$- and $(S,T)$-flow calculation on $T_L(H_{V'})$.
Finally, it follows
\begin{align*}
\omega_H(B_2) \stackrel{\text{Problem \ref{prob:ST}(ii)}}{=} 
\omega_H(B_1) - \omega_{H_{V'}}(B_1) + |f'| \\
\stackrel{\text{Lemma \ref{lemma:ST_subset}}}{\le} \omega_H(B_1) - \omega_{H_{V'}}(B_1) + |f| \\
\stackrel{\text{Problem \ref{prob:ST}(ii)}}{=} \omega_H(B_3) \stackrel{\text{Problem \ref{prob:ST}(i)}}{\le} \omega_H(B_1)
\end{align*}
\end{proof}

We are now able to extract a subhypergraph $H_{V'}$ out of a already bipartitioned hypergraph $H$ and
calculate a minimum $(S,T)$-bipartition of $H_{V'}$ with $S$ and $T$ as defined
in Equation \ref{S_final_border_hyperedges} and \ref{T_final_border_hyperedges}. The resulting
bipartition induced a new cut on $H$ smaller or equal than the old cut. Further, we show with our
modelling technique of $S$ and $T$ that $\Delta_H$ can be calculated with the help of the value 
of a maximum $(S,T)$-flow computation on $T_L(H_{V'})$. Additionaly, we demonstrate that a different
modelling approach of $S$ and $T$ which satisfy both conditions of Problem \ref{prob:ST} can lead to
an improved cut quality of the minimum $(S,T)$-bipartition on the original hypergraph $H$.

\begin{figure}
\centering
\includegraphics[width=0.95\textwidth]{../img/source_and_sink_set/st_modelling_summary.eps}
\caption{Illustration of source and sink set modelling defined in Equation \ref{S_final_border_hyperedges}
         and \ref{T_final_border_hyperedges}. }
\label{img:st_modelling_summary}
\end{figure}

With the given approach we are able to optimize the cut metric of a given
bipartition of a hypergraph $H$. We can transfer those results in order to improve
a $k$-way partition $\Pi = (V_1,\ldots,V_k)$, if the objective is the connectivity
metric. Let $V' \subseteq V_i \cup V_j$ be a subset of the hypernodes of two adjacent
blocks $V_i$ and $V_j$ in the quotient graph. If we optimize the cut of
subhypergraph $H_{V'}$ we simultanously optimize the connectivity metric of $H$.
The reduction of the cut on $H_{V'}$ is then equal with the reduction of
the connectivity on $H$.

\subsection{Most Balanced Minimum Cuts on Hypergraphs}
\label{sec:mbmc_hypergraphs}

Picard and Queyranne \cite{picard1980structure} showed that all minimum $(s,t)$-cuts 
of a graph $G$ are computable with one maximum $(s,t)$-flow computation by 
iterating through all \emph{closed node sets} of the residual graph of $G$. 
The corresponding algorithm is presented in Section \ref{sec:related_mbmc}. \\
We can apply the same algorithm on hypergraphs. A minimum-capacity $(s,t)$-cutset of \ShortT{L}
is equal with a minimum-weight $(s,t)$-cutset of $H$. With the algorithm
of Section \ref{sec:related_mbmc} we are able to find all minimum-capacities
$(s,t)$-cutsets of \ShortT{L}, which are also minimum-weight $(s,t)$-cutsets
of $H$. The corresponding minimum-weight $(s,t)$-bipartitions are all
\emph{closed node sets} of the residual graph of \ShortT{L}. \\
However, when we use e.g. \ShortExT{H}{V'} (see Section \ref{sec:heuer_network})
or \ShortHybrid~(see Section \ref{sec:hybrid_network}) as underlying flow network
some hypernodes are removed from the flow problem. This a problem, if we want to
enumerate all minimum-weight $(s,t)$-bipartitions. The solution for 
this problem is quite simple. After a maximum $(s,t)$-flow calculation
on one of the two mentioned networks we insert all removed hypernodes with
their corresponding edges again into the residual graph of our flow network.
The maximum $(s,t)$-flow is still maximal, otherwise we would had found an \emph{augmenting
path} on the flow network before. We are now able to compute all minimum-weight
$(s,t)$-bipartitions the same way as with \ShortT{L}.

\subsection{A Direct $K$-Way Flow-Based Refinement Framework}
\label{sec:flow_local_search_hypergraph}

We described how a hypergraph $H$ could be transformed into
a flow network \ShortT{L} such that every minimum-capacity $(s,t)$-cutset on \ShortT{L} is a 
minimum-weight  $(s,t)$-cutset on $H$ (see Section \ref{sec:related_lawler}). 
Additionaly, we present techniques to sparsify the
flow network \ShortT{L} \cite{lawler1973} in order to reduce the complexity of 
the flow problem (see Section \ref{sec:opt_flow_network}). 
Further, we show how to configure the source and sink sets on the flow network of a 
subhypergraph $H_{V'}$ (with $V' \subseteq V$) such that a \emph{Max-Flow-Min-Cut} 
computation improves a given bipartition of $H$ (see Section \ref{sec:source_and_sink}). 
Finally, we are able to enumerate all minmum-weight $(s,t)$-cutsets of a subhypergraph 
$H_{V'}$ with one maximum $(s,t)$-flow calculation \cite{picard1980structure}. \\
We will now present our direct $k$-way flow-based refinement framework which we integrated
into the $n$-level hypergraph partitioner \emph{KaHyPar} \cite{heuer2017improving} 
(see Section \ref{sec:kahypar}). Our flow-based refinement approach optimizes
the \emph{connectivity} metric. We used a similiar architecture as proposed
by Sanders and Schulz \cite{sanders2011engineering} (see Section 
\ref{sec:flow_local_search_graph}). The basic concepts of the framework are
illustrated in \autoref{img:flow_framework}. \\
Our maximum flow calculations are embedded into an \emph{Active Block Scheduling}
refinement \cite{holtgrewe2010engineering} (see Section \ref{sec:abs}).
Each time we use flows to improve the connectivity metric of
a given $k$-way partition $\Pi$ we construct the quotient graph $Q$ of $\Pi$. 
Afterwards, we iterate over all edges of $Q$ in random order. For each edge
$(V_i,V_j)$ of $Q$ we grow a flow problem around the cut of the bipartition
induced by $V_i$ and $V_j$. In order to do that we use two \BFS, one only 
touches hypernodes of $V_i$ and the second only touches hypernodes of $V_j$.
The \BFS~are initialized with all hypernodes contained in a cut hyperedge
of the bipartition $(V_i,V_j)$. A pairwise flow-based refinement is embedded
into the \emph{adaptive flow iterations} strategy \cite{sanders2011engineering}
(see Section \ref{sec:adaptive_flow_iterations}) which also determines
the size of the flow problem. \\
After we define the subhypergraph $H_{V'}$, which we use to improve the bipartition
$(V_i,V_j)$ on $H$, we construct one of the flow networks proposed in Section
\ref{sec:opt_flow_network} with sources $S$ and sinks $T$ defined in
Section \ref{sec:source_and_sink}. We implemented two maximum flow algorithms.
One is a slightly modified \emph{augmenting path} algorithm of Emdond \& Karp
\cite{edmonds1972theoretical} (see Section \ref{sec:aug_path}) 
and the second is the \emph{Push-Relabel} algorithm of
Goldberg \& Tarjan \cite{cherkassky1997implementing,goldberg1988new} 
(see Section \ref{sec:push_relabel}). Since we have a 
\emph{Multi-Source-Multi-Sink} problem, we can find several \emph{augmenting paths}
with one \BFS. After we execute a \BFS~on the residual graph, we search 
as many as possible edge-disjoint paths in the resulting \BFS-tree connecting a source $s$
with a sink $t$. Our Goldberg \& Tarjan implementation uses a \emph{FIFO} queue and
the \emph{global relabeling} and \emph{gap} heurisitc \cite{cherkassky1997implementing}.
We do not use an external implementation of a maximum flow algorithm.
Since the \emph{I\textbackslash O} of writing a flow problem to memory and reading the
solution would significant slowdown the performance of our algorithm, because we have
to solve an enormous number of flow problems during the \emph{Active Block Scheduling}
refinement. After determine a maximum $(S,T)$-flow on our flow network we iterate over
all minimum $(S,T)$-bipartitions of $H_{V'}$ \cite{picard1980structure} and choose 
the \emph{Most Balanced Minimum Cut} (see Section \ref{sec:related_mbmc} and 
\ref{sec:mbmc_hypergraphs}) according to our \emph{balanced contraint}. \\
\emph{KaHyPar} is a $n$-level hypergraph partitioner ($|V| = n$) taking the 
multilevel paradigm to its extreme by removing only a single vertex in every level
of the hierarchy \cite{akhremtsev2017engineering} (see Section \ref{sec:kahypar}). 
During the refinement step $n$ local searches are instantiated. Therefore, 
using our flow-based refinement as local search algorithm on each level is not 
applicable, because the performance slowdown would be tremendous. On this
reason we introduce \emph{Flow Execution Policies}. One is to execute our flow-based
refinement on each level $i$ where $i = \beta\cdot j$ with $j \in \mathbb{N}_+$ and
$\beta$ as a predefined tunning parameter. Another approach is to simulate a
multilevel partitioner with $\log(n)$ hierarchies. A flow-based refinement is then
executed on each level $i$ where $i = 2^j$ with $j \in \mathbb{N}_+$. Each policy also
executes the \emph{Active Block Scheduling} refinement on the last level of the
hierarchy. In all remaining levels where no flow is executed, we can use a 
\emph{FM}-based local search algorithm 
\cite{akhremtsev2017engineering,fiduccia1988linear,sanchis1989multiple} (see Section 
\ref{sec:abs}). \\
An observation during the implementation of this framework was that only a minority
of the pairwise refinements based on flows yields to an improvement of the connectivity
metric on a hypergraph $H$. Therefore, we introduce several rules which might prevent
unnecessary flow executions to improve the effectiveness ratio by simultanously speed up
the runtime.

\begin{enumerate}
\item If the cut between two adjacent blocks in the quotient graph is small (e.g. $\le 10$) we
      skip the flow-based refinement on these blocks except on the last level of the hierarchy.
\item If the value of the cut of a minimum $(S,T)$-bipartition on $H_{V'}$ is the same 
      as the cut before, we stop the pairwise refinement on these blocks.
\item If a flow-based refinement did not lead to an improvement on two blocks in all levels
      of the multilevel hierarchy, we only execute flows in the first iteration of 
      \emph{Active Block Scheduling} on these blocks.
\end{enumerate}

\begin{figure}
\centering 
\includegraphics[width=1.0\textwidth]{../img/flow_local_search/flow_framework_hypergraph.eps}
\caption{Illustration of our flow-based refinement framework on hypergraphs.}
\label{img:flow_framework}
\end{figure} 

\newpage

%%%%%%%%%%%%%%%%%%%%%%%%%%%%%%%%%%%%%%%%%%%%%%%%%%%%%%%%%%%%%%%%%%%%%%

\section{Experimental Results}

\newpage

%%%%%%%%%%%%%%%%%%%%%%%%%%%%%%%%%%%%%%%%%%%%%%%%%%%%%%%%%%%%%%%%%%%%%%

\section{Conclusion}

\newpage

%%%%%%%%%%%%%%%%%%%%%%%%%%%%%%%%%%%%%%%%%%%%%%%%%%%%%%%%%%%%%%%%%%%%%%

\bibliographystyle{plain}
\bibliography{../literature}

\end{document}
